% Options for packages loaded elsewhere
\PassOptionsToPackage{unicode}{hyperref}
\PassOptionsToPackage{hyphens}{url}
\PassOptionsToPackage{dvipsnames,svgnames,x11names}{xcolor}
%
\documentclass[
  10pt,
]{article}
\usepackage{amsmath,amssymb}
\usepackage{iftex}
\ifPDFTeX
  \usepackage[T1]{fontenc}
  \usepackage[utf8]{inputenc}
  \usepackage{textcomp} % provide euro and other symbols
\else % if luatex or xetex
  \usepackage{unicode-math} % this also loads fontspec
  \defaultfontfeatures{Scale=MatchLowercase}
  \defaultfontfeatures[\rmfamily]{Ligatures=TeX,Scale=1}
\fi
\usepackage{lmodern}
\ifPDFTeX\else
  % xetex/luatex font selection
  \setmainfont[]{DejaVu Serif}
  \setsansfont[]{DejaVu Sans}
  \setmonofont[]{DejaVu Sans Mono}
  \setmathfont[]{Latin Modern Math}
\fi
% Use upquote if available, for straight quotes in verbatim environments
\IfFileExists{upquote.sty}{\usepackage{upquote}}{}
\IfFileExists{microtype.sty}{% use microtype if available
  \usepackage[]{microtype}
  \UseMicrotypeSet[protrusion]{basicmath} % disable protrusion for tt fonts
}{}
\makeatletter
\@ifundefined{KOMAClassName}{% if non-KOMA class
  \IfFileExists{parskip.sty}{%
    \usepackage{parskip}
  }{% else
    \setlength{\parindent}{0pt}
    \setlength{\parskip}{6pt plus 2pt minus 1pt}}
}{% if KOMA class
  \KOMAoptions{parskip=half}}
\makeatother
\usepackage{xcolor}
\usepackage[margin=2.0cm]{geometry}
\usepackage{longtable,booktabs,array}
\usepackage{calc} % for calculating minipage widths
% Correct order of tables after \paragraph or \subparagraph
\usepackage{etoolbox}
\makeatletter
\patchcmd\longtable{\par}{\if@noskipsec\mbox{}\fi\par}{}{}
\makeatother
% Allow footnotes in longtable head/foot
\IfFileExists{footnotehyper.sty}{\usepackage{footnotehyper}}{\usepackage{footnote}}
\makesavenoteenv{longtable}
\setlength{\emergencystretch}{3em} % prevent overfull lines
\providecommand{\tightlist}{%
  \setlength{\itemsep}{0pt}\setlength{\parskip}{0pt}}
\setcounter{secnumdepth}{-\maxdimen} % remove section numbering
% doc_header.tex -- LaTeX preamble additions used by build_docs.sh for every
% Markdown -> LaTeX -> PDF conversion (supplements, README, submission files).
\usepackage{etoolbox}
% longtable may not be loaded by pandoc when a document has no table
\usepackage{longtable}
% bold vectors: \bm -> \symbf under unicode-math (XeLaTeX), bm package otherwise
\makeatletter
\@ifpackageloaded{unicode-math}{\providecommand{\bm}[1]{\symbf{#1}}}{\usepackage{bm}}
\makeatother
% smaller type in long tables so that wide audit tables fit the text width;
% inside a table a line break is allowed after an underscore of plain text
% (pandoc writes it as \_), so that long identifiers wrap inside their column
% rows are set a little apart, so that symbols with indices do not touch
\AtBeginEnvironment{longtable}{\footnotesize\renewcommand{\arraystretch}{1.45}%
  \let\sptunderscore\_\def\_{\sptunderscore\allowbreak}}
\setlength{\LTcapwidth}{\textwidth}
% allow line breaks after slashes and hyphens inside table cells
\usepackage{xurl}
% code blocks: break over-long lines instead of running into the margin.
% Plain ``` blocks are typeset by pandoc as {verbatim}, highlighted ones
% (```bash) as {Highlighting}; both are re-declared as breakable fancyvrb
% environments.  The continuation line is marked with a small hooked arrow.
\usepackage{fvextra}
\RecustomVerbatimEnvironment{verbatim}{Verbatim}{breaklines,breakanywhere,fontsize=\small}
\makeatletter
\@ifundefined{Highlighting}{}{%
  \RecustomVerbatimEnvironment{Highlighting}{Verbatim}{breaklines,breakanywhere,fontsize=\small,commandchars=\\\{\}}}
\makeatother
% inline code (pandoc writes it as \texttt{...}): never hyphenate -- an inserted
% hyphen could be read as part of a file name -- but allow a line break after
% an underscore, slash, comma, colon or equals sign, so that long paths and
% column lists wrap inside the text width and inside table columns.
% Inside a formula (e.g. the name of a notebook symbol in an equation) the same
% command must first leave math mode, which \text of amsmath does.
\ExplSyntaxOn
\tl_new:N \l__spt_code_tl
\cs_new_protected:Npn \__spt_texttt:n #1
  {
    \tl_set:Nn \l__spt_code_tl {#1}
    \regex_replace_all:nnN { ([\/\,\:\=]) } { \1 \c{allowbreak} } \l__spt_code_tl
    \regex_replace_all:nnN { \c{_} } { \c{_} \c{allowbreak} } \l__spt_code_tl
    \group_begin: \ttfamily \hyphenchar \font = -1 \scan_stop: \l__spt_code_tl \group_end:
  }
\RenewDocumentCommand \texttt { m }
  {
    \mode_if_math:TF { \text { \__spt_texttt:n {#1} } } { \__spt_texttt:n {#1} }
  }
\ExplSyntaxOff
% A figure is written in the Markdown source as an image paragraph followed by
% a caption paragraph ("**Figure S1.** ..."), so that the caption is also
% visible on GitHub.  An image that starts a paragraph is centered.
\usepackage{graphicx}% pandoc loads it only for documents that contain an image
\NewCommandCopy\sptincludegraphics\includegraphics
\RenewDocumentCommand\includegraphics{O{} m}{%
  \ifvmode\begin{center}\sptincludegraphics[#1]{#2}\end{center}%
  \else\sptincludegraphics[#1]{#2}\fi}
% keep figures where they are referenced
\usepackage{float}
\floatplacement{figure}{H}
% running header with the document title
\usepackage{fancyhdr}
\pagestyle{fancy}
\fancyhf{}
\fancyhead[L]{\small PROJECT spacetimes --- documentation and provenance}
\fancyhead[R]{\small\thepage}
\renewcommand{\headrulewidth}{0.4pt}
% 16x16 matrices: tighter columns so that a full matrix fits the text width
\setlength{\arraycolsep}{3.2pt}
% The documents of this project contain long tagged symbols, file paths and
% identifiers that cannot be hyphenated.  A paragraph is therefore allowed to
% have looser lines before TeX gives up and lets a line stick out of the margin.
\AtBeginDocument{\tolerance=2000 \emergencystretch=5em \hbadness=2000}
\ifLuaTeX
  \usepackage{selnolig}  % disable illegal ligatures
\fi
\IfFileExists{bookmark.sty}{\usepackage{bookmark}}{\usepackage{hyperref}}
\IfFileExists{xurl.sty}{\usepackage{xurl}}{} % add URL line breaks if available
\urlstyle{same}
\hypersetup{
  pdftitle={P1 --- MASTER provenance: the symbols of PROJECT spacetimes},
  pdfauthor={Prepared for Patrick Lee Nash, Ph.D.},
  colorlinks=true,
  linkcolor={blue},
  filecolor={Maroon},
  citecolor={blue},
  urlcolor={blue},
  pdfcreator={LaTeX via pandoc}}

\title{P1 --- MASTER provenance: the symbols of PROJECT spacetimes}
\usepackage{etoolbox}
\makeatletter
\providecommand{\subtitle}[1]{% add subtitle to \maketitle
  \apptocmd{\@title}{\par {\large #1 \par}}{}{}
}
\makeatother
\subtitle{Which symbol denotes which object in every original file,
where the files contradict one another, and the unique symbols that the
project uses instead}
\author{Prepared for Patrick Lee Nash, Ph.D.}
\date{November 12, 2026}

\begin{document}
\maketitle

\hypertarget{purpose-and-the-result-in-one-page}{%
\section{0. Purpose, and the result in one
page}\label{purpose-and-the-result-in-one-page}}

Hard Rule 5 of the task specification says that the original files of
this project often use the \emph{same} symbol for completely
\emph{different} objects, that this problem must be handled completely
and without ambiguity, and that the proof of the solution is a MASTER
provenance: a clear account of the symbol identifications and of the
remedies, for each and every relevant file, followed by the consistent
use of unique symbols throughout the project. This document is that
provenance.

\textbf{The original files} are

\begin{itemize}
\tightlist
\item
  the 14 Mathematica notebooks of the project folder (identifiers N01 to
  N14; Sec. 5.1),
\item
  nine published papers (D01 to D09; Sec. 5.2), and
\item
  the files of the author's earlier project NCSS, which the task
  specification asks to be read for context (Sec. 5.3).
\end{itemize}

The task specification calls the original files ``yours and mine'': some
were written by the author, some were generated for him (the notebook
\texttt{7+1.nb}, for example, says of itself that a script produced it).
This document makes no difference between them.

\textbf{The problem, measured.} A program written for this purpose finds
1921 different symbols that the notebooks define. 350 of them are
defined in two or more notebooks, and for 76 of these the defining
statements differ from notebook to notebook (Appendix A lists all of
them). The papers and the project NCSS add collisions of their own: the
index of the \(8\times8\) matrices \(\tau^A\) runs from 1 to 8 in the
papers D07 and D08 and from 0 to 7 in the notebooks; the ``strength of
the force'' \(D\) of NCSS is one half of the quantity that carries this
letter here.

\textbf{The remedy} (Sec. 2) is one rule with three parts.

\begin{enumerate}
\def\labelenumi{\arabic{enumi}.}
\tightlist
\item
  Every object that the project uses has exactly one symbol. The symbols
  are listed in a registry, \texttt{provenance/symbols.csv} (113
  entries; Sec. 3).
\item
  A symbol of the project cannot be mistaken for a symbol of an original
  file. Every name that the programs and data files use for a matrix or
  a field carries the tag \textbf{1644} (4+4 spacetime) or \textbf{1671}
  (7+1 spacetime), and every name that belongs to the new scalar field
  carries the suffix \textbf{VS}. No symbol of any notebook contains
  such a tag or suffix.
\item
  An object taken from an original file is identified once, here and in
  the document that introduces it, by \emph{file and original symbol};
  after that only the symbol of the project is used.
\end{enumerate}

\textbf{The proof} (Sec. 8) is a set of 9 machine checks, all of which
pass. They establish that the registry has no duplicate, that none of
its 78 names is a symbol of any notebook, that every tagged name
occurring in any data file or document of the project is registered,
that the manuscripts contain none of the original names of the objects
that were taken over, and that every plain letter of the registry that
coincides with a symbol of a notebook is listed with both meanings.

\textbf{What a reader should take away.} If a symbol carries a tag or
the suffix VS, look it up in Sec. 3: it has one meaning, in every file
of the project. If a formula of an \emph{original} file is to be
compared with the project, look up the file in Sec. 5 and the symbol in
Sec. 6: the same letter may denote something else there.

\hypertarget{how-to-reproduce-the-lists-and-the-checks}{%
\subsection{0.1 How to reproduce the lists and the
checks}\label{how-to-reproduce-the-lists-and-the-checks}}

From the root of the repository, with the notebooks in that folder (they
are not part of the repository; \texttt{provenance/inputs\_sha256.txt}
pins them, and the first script refuses to run if a checksum differs):

\begin{verbatim}
python3 tools/scan_notebook_symbols.py          # symbols of the notebooks -> provenance/generated/
python3 code/scripts/write_symbol_registry.py   # the registry             -> provenance/symbols.csv
python3 code/scripts/check_symbols.py           # the checks               -> data/symbol_checks.txt
python3 code/scripts/symbol_concordance.py      # letters in the documents -> data/symbol_concordance.csv
python3 code/scripts/assemble_docs.py docs_src/P1_master_symbols.src.md
bash code/scripts/build_docs.sh provenance/P1_master_symbols.md
\end{verbatim}

The scripts need only Python 3. The first converts each notebook to text
with \texttt{tools/nb2txt.py} and takes about fifteen seconds; what it
counts as ``a notebook defines a symbol'' is stated in its header and in
Appendix A.

\begin{center}\rule{0.5\linewidth}{0.5pt}\end{center}

\hypertarget{the-problem}{%
\section{1. The problem}\label{the-problem}}

Four examples show the kinds of collision that occur. Each is treated in
Secs. 5 and 6.

\emph{The same name for different matrices.} In \texttt{4+4.nb} (N01)
the name \texttt{γ} is a list of sixteen real \(4\times4\) matrices, a
basis of all real \(4\times4\) matrices. In
\texttt{Dirac\_Nash\_Identity.nb} (N05) \texttt{γ} is a list of four
\(4\times4\) matrices. In
\texttt{Generalized\ \_Kronecker\_Delta\_4+4.nb} (N09) \texttt{γ} is a
tensor index. In the paper D06, \(\gamma^{AB}\) are fifteen real
\(4\times4\) matrices labelled by a pair of indices with six values. In
this project \(\gamma\) denotes the eight real \(16\times16\) Dirac
matrices of an eight-dimensional spacetime, and there are two different
sets of them.

\emph{The same name for different coordinates.} Nine notebooks define a
list \texttt{X} of coordinates. In most of them it is
\(\{x0,\dots,x7\}\) with \(x0\) a hidden space-like coordinate and
\(x4\) the time. In \texttt{7+1.nb} (N02) the same eight names are used,
but \(x0\) is the time. In \texttt{Dirac\_Nash\_Identity.nb} (N05)
\texttt{X} has six entries, \(x1\) to \(x6\).
\texttt{Einstein-Lovelock-Eternal.nb} (N08) uses the coordinates \(x1\)
to \(x8\), and there \(x8\) plays the part that \(x0\) plays elsewhere.

\emph{Two objects under almost the same name in one file.}
\texttt{7+1.nb} (N02) contains two sets of Dirac matrices for the 7+1
signature, the real set \texttt{Γ17} and the complex set \texttt{Γ71};
two metrics, \texttt{η17} and \texttt{η71} \(=-\)\texttt{η17}; and two
chirality matrices, \texttt{χ17} and \texttt{χ71}. A reader who
exchanges the two digits obtains a different theory.

\emph{The same letter for quantities that differ by a factor.} The
project NCSS calls \(D=\alpha m_\nu/2m_\chi\) the strength of the force
between two fermions; two fermions bind for \(D>0.83990\). Here the
corresponding number is \(\alpha_{VS}m/m_s\), without the factor
\(\tfrac12\), and two fermions bind when it exceeds \(1.67981\). A
reader who carries a value of \(D\) from one project to the other is
wrong by a factor of two.

The project itself made errors of this kind before the rule of Sec. 2
was enforced by a program. They are recorded in Sec. 7 and in the error
log (\texttt{provenance/P9\_error\_log}, entries E7, E35, E41 and E42).

\begin{center}\rule{0.5\linewidth}{0.5pt}\end{center}

\hypertarget{the-remedy}{%
\section{2. The remedy}\label{the-remedy}}

\textbf{Rule 1: one object, one symbol; one symbol, one object.} The
registry of Sec. 3 is the authority. A new object gets a new entry
before it is used in a document.

\textbf{Rule 2: tags.} The tag \textbf{1644} marks an object of the 4+4
spacetime, the spacetime of the field dirac1644 (16 components;
signature 4+4). The tag \textbf{1671} marks an object of the 7+1
spacetime, the spacetime of the field dirac1671. In formulas the tag is
a subscript, \(\gamma_{1644}^A\); in programs and data files it is the
end of the name, \texttt{gamma1644}. The suffix \textbf{VS} marks the
quantities that define the Vorton-Scalar: \(\Phi_{VS}\), \texttt{PhiVS}.

\textbf{Rule 3: identification by file and original symbol.} When an
object is taken from an original file, the document that introduces it
states the file and the original symbol, for example
``\(\gamma_{1644}^A\) is \texttt{T16\^{}A{[}A{]}} of \texttt{4+4.nb}'',
and a machine check compares the two entry by entry where that is
possible (checks \texttt{44-N1} to \texttt{44-N3}, \texttt{44-P3},
\texttt{71-N1} to \texttt{71-N3} of the program
\texttt{spacetimes\ matrices}, which read the matrices that Mathematica
printed in the notebooks). The original symbol is not used again.

\textbf{Rule 4: the scope of a dropped tag.} Inside a section that deals
with one spacetime only, the tag may be dropped in derivations; the
equations that define an object or state a result carry it. In the
sections that the two manuscripts share, \(\beta\), \(\alpha^j\),
\(\chi\), \(\mathsf h\) and \(m\) denote the matrices and the mass of
the spacetime at hand (Sec. 3.4). Two symbols of this kind deserve a
warning, because they are the same letter for corresponding but
different matrices: in the 7+1 spacetime
\(\chi\equiv\chi_{1671}=i\,\omega_{1671}\) is complex, in the 4+4
spacetime \(\chi\equiv\omega_{1644}\) is real; and the adjoint
\(\bar\psi\equiv\psi^\dagger\beta\) of the quantized theories differs,
in 4+4, from the covariant adjoint
\(\tilde\psi\equiv\psi^\dagger C_{1644}^{-}\).

\textbf{Rule 5: letters without a tag.} A document of physics cannot
give every quantity a letter of its own; the alphabet is too small.
Letters that the project uses for more than one quantity are listed in
Sec. 4, with the subscript, argument or section that tells the meanings
apart. Plain letters of the project that coincide with a symbol of a
notebook are listed, with both meanings, in Sec. 6.3; a machine check
(\texttt{S9}) guarantees that this list is complete for the symbols of
the registry.

\begin{center}\rule{0.5\linewidth}{0.5pt}\end{center}

\hypertarget{the-registry-the-symbols-of-the-project}{%
\section{3. The registry: the symbols of the
project}\label{the-registry-the-symbols-of-the-project}}

The tables of this section are generated from
\texttt{provenance/symbols.csv}. The column ``Name in files'' gives the
names used in the output of the programs and in the data files;
``Defined in'' uses the abbreviations DM for the provenance ``dirac
matrices'', S1 to S4 for the supplements and M for the two manuscripts
(the section numbers of the manuscripts are the same in both).
``Origin'' names the notebook (N01 to N14) or the paper (D01 to D09) and
the original symbol; ``new'' means that the project introduced the
object. The word NOT marks an original symbol that is easily confused
with the right one.

\hypertarget{coordinates-and-indices}{%
\subsection{3.1 Coordinates and indices}\label{coordinates-and-indices}}

\begin{longtable}[]{@{}
  >{\raggedright\arraybackslash}p{(\columnwidth - 8\tabcolsep) * \real{0.1446}}
  >{\raggedright\arraybackslash}p{(\columnwidth - 8\tabcolsep) * \real{0.1325}}
  >{\raggedright\arraybackslash}p{(\columnwidth - 8\tabcolsep) * \real{0.3976}}
  >{\raggedright\arraybackslash}p{(\columnwidth - 8\tabcolsep) * \real{0.1205}}
  >{\raggedright\arraybackslash}p{(\columnwidth - 8\tabcolsep) * \real{0.2048}}@{}}
\toprule\noalign{}
\begin{minipage}[b]{\linewidth}\raggedright
Symbol
\end{minipage} & \begin{minipage}[b]{\linewidth}\raggedright
Name in files
\end{minipage} & \begin{minipage}[b]{\linewidth}\raggedright
Meaning
\end{minipage} & \begin{minipage}[b]{\linewidth}\raggedright
Defined in
\end{minipage} & \begin{minipage}[b]{\linewidth}\raggedright
Origin
\end{minipage} \\
\midrule\noalign{}
\endhead
\bottomrule\noalign{}
\endlastfoot
\(X=\{x0,\dots,x7\}\) & \texttt{x0} \texttt{x1} \texttt{x2} \texttt{x3}
\texttt{x4} \texttt{x5} \texttt{x6} \texttt{x7} & the eight coordinates
of both spacetimes, never renamed: \(x0\) hidden space-like;
\(x1,x2,x3\) ordinary space; \(x4\) ordinary time; \(x5,x6,x7\)
additional coordinates (time-like in 4+4, space-like in 7+1) & DM 0.1; M
I & N01 (text at its head) and N04 \texttt{X} (section DEFCHART): the
same eight names with the same meaning. N02 uses the same eight names
with \texttt{x0} as the time: related by swap04 \\
\(\mathrm{swap}_{04}\) & \texttt{swap04} & the exchange of the index
values 0 and 4; it carries the labels of \texttt{7+1.nb} (time at index
0) into those of the project (time at index 4) & DM 0.1, 2.1 & new \\
\(A,B,C,D\) & --- & index with the values \(0,\dots,7\); always labels
the coordinate \(xA\) & DM 0.1; M I & new \\
\(a,b\) & --- & in S1 to S4 and in the manuscripts: index of the four
extra axes, values 0, 5, 6, 7. In DM Sec. 1.8.4 and in the texts of the
checks \texttt{44-D6}, \texttt{44-D7}: the four space-like axes of 4+4,
values 0, 1, 2, 3 (stated there) & M I; S1 0.1; DM 1.8.4 & new \\
\(t\) (index) & --- & index of the three additional coordinates, values
5, 6, 7 & DM 1.8.4; M I & new \\
\(i,j\) & --- & index of space-like coordinates; values 1, 2, 3
(ordinary space) unless a sum is said to run over all space-like
coordinates of the spacetime & M I & new \\
\(h\) (index) & --- & index with the values 1, 2, 3 that labels the
\(4\times4\) matrices \(s_h\), \(t_h\) & DM 1.1; M II & N01, N02, N05
\texttt{h} in \texttt{s4by4{[}h{]}}, \texttt{t4by4{[}h{]}} \\
\end{longtable}

\hypertarget{matrices-of-the-44-spacetime}{%
\subsection{3.2 Matrices of the 4+4
spacetime}\label{matrices-of-the-44-spacetime}}

\begin{longtable}[]{@{}
  >{\raggedright\arraybackslash}p{(\columnwidth - 8\tabcolsep) * \real{0.1446}}
  >{\raggedright\arraybackslash}p{(\columnwidth - 8\tabcolsep) * \real{0.1325}}
  >{\raggedright\arraybackslash}p{(\columnwidth - 8\tabcolsep) * \real{0.3976}}
  >{\raggedright\arraybackslash}p{(\columnwidth - 8\tabcolsep) * \real{0.1205}}
  >{\raggedright\arraybackslash}p{(\columnwidth - 8\tabcolsep) * \real{0.2048}}@{}}
\toprule\noalign{}
\begin{minipage}[b]{\linewidth}\raggedright
Symbol
\end{minipage} & \begin{minipage}[b]{\linewidth}\raggedright
Name in files
\end{minipage} & \begin{minipage}[b]{\linewidth}\raggedright
Meaning
\end{minipage} & \begin{minipage}[b]{\linewidth}\raggedright
Defined in
\end{minipage} & \begin{minipage}[b]{\linewidth}\raggedright
Origin
\end{minipage} \\
\midrule\noalign{}
\endhead
\bottomrule\noalign{}
\endlastfoot
\(\eta_{1644}\) & \texttt{eta1644} & metric of the 4+4 spacetime,
\(\mathrm{diag}(+1,+1,+1,+1,-1,-1,-1,-1)\): space-like \(=+1\) & DM 1.4,
Eq. (1.3) & N01, N02, N05, N10--N12 \texttt{η4488}; in N01 also
\texttt{η8} and \texttt{g\_AB}; D08 \(G_{AB}\) with the index value 8 in
the place of the value 0 \\
\(s_h\), \(t_h\) & --- & the three self-dual and the three
anti-self-dual real antisymmetric \(4\times4\) matrices & DM 1.1; M II
(4+4), App. A (7+1) & N01, N02, N05, N10--N12 \texttt{s4by4{[}h{]}},
\texttt{t4by4{[}h{]}}; the letters are those of D06 (\(s^h\) self-dual,
\(t^h\) anti-self-dual) \\
\(\tau_{1644}^A\), \(\bar\tau_{1644}^A\) & \texttt{tau1644}
\texttt{taubar1644} & the real \(8\times8\) blocks of the 4+4 Dirac
matrices (reduced generators of the split octonion algebra) & DM 1.1,
App. B; M II & N01, N10--N12 \texttt{τ{[}A{]}} and \texttt{τ} with an
overbar; N02 \texttt{τ{[}A{]}}, \texttt{τbar{[}A{]}}; matrices of this
kind are the \(\tau^A\), \(\bar\tau^A\) of D07 and D08, whose index runs
from 1 to 8 (the value 8 there is the value 0 here) \\
\(\sigma_{88}\) & \texttt{sigma88} & the \(8\times8\) matrix with the
\(4\times4\) blocks \((0,1_4;\,1_4,0)\) & DM 1.1 & N01 \texttt{σ},
\texttt{σ88}; N02 \texttt{σ}; N05 \texttt{σ88} \\
\(\gamma_{1644}^A\) & \texttt{gamma1644} & the eight real \(16\times16\)
Dirac matrices of the 4+4 spacetime & DM 1.1, 1.2; M II & N01, N10--N12
\texttt{T16\^{}A{[}A{]}}; N02 \texttt{T16{[}A{]}}; the \(t^A\) of D08
are built in the same way \\
\(\Sigma_{1644}^{AB}\) & \texttt{Sigma1644} &
\(\tfrac14[\gamma_{1644}^A,\gamma_{1644}^B]\), the scaled commutators;
they span \(\mathfrak{so}(4,4)\) & DM 1.5; M II & new name; the same
matrices are \texttt{SAB} and \texttt{sAB{[}A{]}{[}B{]}} of N10--N12;
NOT \texttt{sAB} of N02, which is a sign \\
\(\omega_{1644}\) & \texttt{omega1644} & the product
\(\gamma_{1644}^0\cdots\gamma_{1644}^7\); real chirality matrix of 4+4,
\(\mathrm{diag}(-1_8,+1_8)\) & DM 1.6; M II & N01
\texttt{T16\^{}A{[}8{]}}; N02 \texttt{T16{[}8{]}} \\
\(P_{1644}^{L}\), \(P_{1644}^{R}\) & \texttt{PL1644} \texttt{PR1644} &
the real chiral projectors \(\tfrac12(1\mp\omega_{1644})\) & DM 1.6; M
II & N01 \texttt{P\_L}, \texttt{P\_R}; N02 \texttt{PL44},
\texttt{PR44} \\
\(C_{1644}^{-}\), \(C_{1644}^{+}\) & \texttt{Cm1644} \texttt{Cp1644} &
the products \(\gamma^0\gamma^1\gamma^2\gamma^3\) and
\(\gamma^4\gamma^5\gamma^6\gamma^7\) of the 4+4 matrices;
\(C^{\mp}\gamma^A(C^{\mp})^{-1}=\mp(\gamma^A)^T\) & DM 1.7; M II & N01,
N02, N10--N12 \texttt{σ16} (\(C^{-}\)); N02 \texttt{S44} (\(C^{+}\)) \\
\(W_{1644}\) & \texttt{W1644} & \(i\,C_{1644}^{-}\gamma_{1644}^4\), the
grading: Hermitian, \(W^2=1\); density of the covariant current & DM
1.8.4; M III & new \\
\(K_{1644}\) & \texttt{K1644} & generator of the evolution of dirac1644
in \(x4\): \(\partial_4\psi=K\psi\) & DM 1.8.4, Eq. (1.15); M III &
new \\
\(\beta_{1644}\), \(\alpha_{1644}^j\) & \texttt{beta1644}
\texttt{alpha1644} & the Hermitian matrices \(-i\gamma_{1644}^4\) and
\(-\gamma_{1644}^4\gamma_{1644}^j\) (\(j=0,1,2,3\)) of the quantized 4+4
theory & S1 2; M III & new \\
\(\mathsf h_{1644}\) & \texttt{h1644} & Hamiltonian matrix of the
quantized 4+4 theory, \(iK_{1644}\) for fields independent of
\(x5,x6,x7\) & DM 1.8.4; M III & new \\
\(\Pi_{1644}\) & \texttt{Pi1644} & the projector
\((i\gamma_{1644}^Ak_A+m_{1644})/2m_{1644}\) on the plane-wave solutions
of the worked example & DM 1.8.1; M III & new \\
\end{longtable}

\hypertarget{matrices-of-the-71-spacetime}{%
\subsection{3.3 Matrices of the 7+1
spacetime}\label{matrices-of-the-71-spacetime}}

\begin{longtable}[]{@{}
  >{\raggedright\arraybackslash}p{(\columnwidth - 8\tabcolsep) * \real{0.1446}}
  >{\raggedright\arraybackslash}p{(\columnwidth - 8\tabcolsep) * \real{0.1325}}
  >{\raggedright\arraybackslash}p{(\columnwidth - 8\tabcolsep) * \real{0.3976}}
  >{\raggedright\arraybackslash}p{(\columnwidth - 8\tabcolsep) * \real{0.1205}}
  >{\raggedright\arraybackslash}p{(\columnwidth - 8\tabcolsep) * \real{0.2048}}@{}}
\toprule\noalign{}
\begin{minipage}[b]{\linewidth}\raggedright
Symbol
\end{minipage} & \begin{minipage}[b]{\linewidth}\raggedright
Name in files
\end{minipage} & \begin{minipage}[b]{\linewidth}\raggedright
Meaning
\end{minipage} & \begin{minipage}[b]{\linewidth}\raggedright
Defined in
\end{minipage} & \begin{minipage}[b]{\linewidth}\raggedright
Origin
\end{minipage} \\
\midrule\noalign{}
\endhead
\bottomrule\noalign{}
\endlastfoot
\(\eta_{1671}\) & \texttt{eta1671} & metric of the 7+1 spacetime,
\(\mathrm{diag}(-1,-1,-1,-1,+1,-1,-1,-1)\): time-like \(=+1\), time at
index 4 & DM 2.4, Eq. (2.3); M II & N02 \texttt{η17} with the index
values 0 and 4 exchanged; NOT \texttt{η71} \\
\(\tau_{1671}^A\) & \texttt{tau1671} & the real \(8\times8\) matrices of
the 7+1 construction: \(\tau_{1671}^4=1_8\), the other seven represent
the imaginary octonion units & DM 2.1, App. B; M II & N02
\texttt{τO{[}swap04(A){]}} \\
\(\gamma_{1671}^A\) & \texttt{gamma1671} & the eight real \(16\times16\)
Dirac matrices of the 7+1 spacetime (contravariant index) & DM 2.1, 2.2;
M II & N02 \texttt{Γ17up{[}swap04(A){]}}; NOT \texttt{Γ71{[}A{]}} \\
\(\Sigma_{1671}^{AB}\) & \texttt{Sigma1671} &
\(\tfrac14[\gamma_{1671}^A,\gamma_{1671}^B]\), the scaled commutators;
they span \(\mathfrak{so}(7,1)\) & DM 2.5; M II & N02
\texttt{Σ17{[}A,B{]}}, built from the \texttt{Γ17} with lower index:
\(\Sigma_{1671}^{AB}=\eta_{1671}^{AA}\eta_{1671}^{BB}\,\)\texttt{Σ17{[}A\textquotesingle{},B\textquotesingle{}{]}}
with \(A'=\mathrm{swap}_{04}(A)\), \(B'=\mathrm{swap}_{04}(B)\) \\
\(\omega_{1671}\) & \texttt{omega1671} & the product
\(\gamma_{1671}^0\cdots\gamma_{1671}^7\); real, \(\omega_{1671}^2=-1\) &
DM 2.6; M II & N02 \texttt{Γ17{[}8{]}}, \texttt{ω17} \\
\(\chi_{1671}\) & \texttt{chi1671} & \(i\,\omega_{1671}\), the (complex)
chirality matrix of 7+1 & DM 2.6; M II & N02 \texttt{χ17}; NOT
\texttt{χ71} \\
\(P_{1671}^{L}\), \(P_{1671}^{R}\) & \texttt{PL1671} \texttt{PR1671} &
the complex chiral projectors \(\tfrac12(1\mp\chi_{1671})\) & DM 2.6; M
II & N02 \texttt{PL17}, \texttt{PR17}; NOT \texttt{PL71},
\texttt{PR71} \\
\(U_{1671}^{c}\) & \texttt{Uc1671} & the unitary matrix
\((1-i\gamma_{1671}^4)/\sqrt2\) that carries the complex projectors of
7+1 into the real projectors of 4+4 & DM 2.6, Eq. (2.5); M II & new \\
\(C_{1671}^{+}\), \(C_{1671}^{-}\) & \texttt{Cp1671} \texttt{Cm1671} &
\(\gamma_{1671}^4\) and \(\gamma_{1671}^4\omega_{1671}\);
\(C^{\pm}\gamma^A(C^{\pm})^{-1}=\pm(\gamma^A)^T\) & DM 2.7; M II &
new \\
\(\beta_{1671}\), \(\alpha_{1671}^j\) & \texttt{beta1671}
\texttt{alpha1671} & the real symmetric matrices \(\gamma_{1671}^4\) and
\(\gamma_{1671}^4\gamma_{1671}^j\) (seven space indices \(j\)) of the
Hamiltonian & DM 2.8.2; M II & N02 \texttt{βD},
\texttt{αD{[}swap04(j){]}} \\
\(\mathsf h_{1671}\) & \texttt{h1671} & Hamiltonian matrix of the 7+1
theory, \(\alpha_{1671}^jp^j+\beta_{1671}m_{1671}\) & DM 2.8.2; M III &
N02 \texttt{h} \\
\(\Pi_{1671}\) & \texttt{Pi1671} & the projector
\((\gamma_{1671}^Ak_A+m_{1671})/2m_{1671}\) of the worked example & DM
2.8.1 & new \\
\(E_{1671}^{L}\), \(E_{1671}^{R}\) & \texttt{EL1671} \texttt{ER1671} &
the chiral frames: \(16\times8\) matrices whose columns span the two
eigenspaces of \(\chi_{1671}\) & DM 2.8.3; M III & N02 \texttt{EL},
\texttt{ER} \\
\(M_{1671}\), \(U_{1671}^{+}\), \(V_{1671}^{-}\) & \texttt{M1671}
\texttt{Upos1671} \texttt{Vneg1671} & the \(8\times8\) matrix
\(m_{1671}+p^j\tau_{1671}^j\) and the matrices of positive- and
negative-energy spinors of the worked example (check \texttt{71-D8}) &
DM 2.8.2 & N02 \texttt{M} of its Sec. 7.4 (in N02 the letter \texttt{M}
also denotes other matrices and a mass); the two spinor matrices are new
names \\
\end{longtable}

The two sets of Dirac matrices are related exactly (check
\texttt{71-B4}): \(\gamma_{1671}^4=\gamma_{1644}^0\);
\(\gamma_{1671}^j=-\gamma_{1644}^0\gamma_{1644}^j\) for \(j=1,2,3\);
\(\gamma_{1671}^0=+\gamma_{1644}^0\gamma_{1644}^5\gamma_{1644}^6\gamma_{1644}^7\),
\(\gamma_{1671}^5=-\gamma_{1644}^0\gamma_{1644}^4\gamma_{1644}^6\gamma_{1644}^7\),
\(\gamma_{1671}^6=+\gamma_{1644}^0\gamma_{1644}^4\gamma_{1644}^5\gamma_{1644}^7\),
\(\gamma_{1671}^7=-\gamma_{1644}^0\gamma_{1644}^4\gamma_{1644}^5\gamma_{1644}^6\).
A formula that is true for one set is therefore in general false for the
other, and the tag is the only thing that tells them apart.

\hypertarget{fields-and-their-equations}{%
\subsection{3.4 Fields and their
equations}\label{fields-and-their-equations}}

\begin{longtable}[]{@{}
  >{\raggedright\arraybackslash}p{(\columnwidth - 8\tabcolsep) * \real{0.1446}}
  >{\raggedright\arraybackslash}p{(\columnwidth - 8\tabcolsep) * \real{0.1325}}
  >{\raggedright\arraybackslash}p{(\columnwidth - 8\tabcolsep) * \real{0.3976}}
  >{\raggedright\arraybackslash}p{(\columnwidth - 8\tabcolsep) * \real{0.1205}}
  >{\raggedright\arraybackslash}p{(\columnwidth - 8\tabcolsep) * \real{0.2048}}@{}}
\toprule\noalign{}
\begin{minipage}[b]{\linewidth}\raggedright
Symbol
\end{minipage} & \begin{minipage}[b]{\linewidth}\raggedright
Name in files
\end{minipage} & \begin{minipage}[b]{\linewidth}\raggedright
Meaning
\end{minipage} & \begin{minipage}[b]{\linewidth}\raggedright
Defined in
\end{minipage} & \begin{minipage}[b]{\linewidth}\raggedright
Origin
\end{minipage} \\
\midrule\noalign{}
\endhead
\bottomrule\noalign{}
\endlastfoot
\(\psi_{1644}\) & \texttt{dirac1644} & the field dirac1644: 16 complex
Grassmann components on the 4+4 spacetime;
\((\gamma_{1644}^A\partial_A-m_{1644})\psi_{1644}=0\) & DM 1.8; M III &
new; \texttt{Ψ16} of N01 and N10--N12 is a column of 16 real quantities
(in N10--N12 functions of \texttt{x0}, \texttt{x4} called the wave
function of the universe), not this field \\
\(\psi_{1644}^{L}\), \(\psi_{1644}^{R}\) & \texttt{psiL1644}
\texttt{psiR1644} & the two 8-component chiral parts of dirac1644 (upper
and lower eight components) & DM 1.8.5; M III & new \\
\(m_{1644}\) & \texttt{m1644} & mass of the field dirac1644 & DM 1.8.1 &
new \\
\(\tilde\psi\) & --- & the covariant adjoint
\(\psi^\dagger C_{1644}^{-}\) of the 4+4 theory (invariant under
\(\mathrm{Spin}_0(4,4)\)) & DM 1.8.2; M III (4+4) & new \\
\(j_{1644}^A\) & \texttt{j1644} & the conserved current
\(i\,\tilde\psi\gamma_{1644}^A\psi\) of the covariant 4+4 theory
(fermion number) & DM 1.8.2, Eq. (1.13); M III & new \\
\(\psi_{1671}\) & \texttt{dirac1671} & the field dirac1671: 16 complex
Grassmann components on the 7+1 spacetime;
\((i\gamma_{1671}^A\partial_A-m_{1671})\psi_{1671}=0\) & DM 2.8; M III &
N02 \texttt{dirac16complex} \\
\(\phi_{1671}^{L}\), \(\phi_{1671}^{R}\) & \texttt{phiL1671}
\texttt{phiR1671} & the two 8-component chiral parts of dirac1671 & DM
2.8.3; M III & N02 \texttt{φL}, \texttt{φR} \\
\(m_{1671}\) & \texttt{m1671} & mass of the field dirac1671 & DM 2.8.1 &
N02 \texttt{m}, \texttt{mD} \\
\(\bar\psi\) & --- & the adjoint \(\psi^\dagger\beta\) of the quantized
theories (\(\beta=\beta_{1671}\) or \(\beta_{1644}\)); in 4+4 it differs
from \(\tilde\psi\) & S1 1.2, 3.2; M V & new \\
\(\beta\), \(\alpha^j\), \(\chi\), \(\mathsf h\), \(m\) & --- & in the
sections common to both spacetimes: the matrices and the mass of the
spacetime at hand, i.e.~the tagged symbols with the tag dropped;
\(\chi\equiv\chi_{1671}\) in 7+1 and \(\chi\equiv\omega_{1644}\) in 4+4
& S1 2, Eq. (2.1); M V (4+4), Eq. for \(\chi\) & new \\
\(U(\vec p)\) & --- & the \(16\times8\) matrix of orthonormal
positive-energy spinors at momentum \(\vec p\) & S1 4.1; M VI & new \\
\end{longtable}

\hypertarget{symmetries-pairs-and-the-order-parameter}{%
\subsection{3.5 Symmetries, pairs and the order
parameter}\label{symmetries-pairs-and-the-order-parameter}}

\begin{longtable}[]{@{}
  >{\raggedright\arraybackslash}p{(\columnwidth - 8\tabcolsep) * \real{0.1446}}
  >{\raggedright\arraybackslash}p{(\columnwidth - 8\tabcolsep) * \real{0.1325}}
  >{\raggedright\arraybackslash}p{(\columnwidth - 8\tabcolsep) * \real{0.3976}}
  >{\raggedright\arraybackslash}p{(\columnwidth - 8\tabcolsep) * \real{0.1205}}
  >{\raggedright\arraybackslash}p{(\columnwidth - 8\tabcolsep) * \real{0.2048}}@{}}
\toprule\noalign{}
\begin{minipage}[b]{\linewidth}\raggedright
Symbol
\end{minipage} & \begin{minipage}[b]{\linewidth}\raggedright
Name in files
\end{minipage} & \begin{minipage}[b]{\linewidth}\raggedright
Meaning
\end{minipage} & \begin{minipage}[b]{\linewidth}\raggedright
Defined in
\end{minipage} & \begin{minipage}[b]{\linewidth}\raggedright
Origin
\end{minipage} \\
\midrule\noalign{}
\endhead
\bottomrule\noalign{}
\endlastfoot
\(U(1)_F\) & --- & the phase symmetry \(\psi\to e^{i\epsilon}\psi\):
fermion number & S1 1.2; M V & new \\
\(U(1)_{VS}\) & \texttt{UVS} & the chiral symmetry
\(\psi\to e^{i\epsilon\chi}\psi\),
\(\Phi_{VS}\to e^{-2i\epsilon}\Phi_{VS}\), broken by the vacuum value of
the Vorton-Scalar & S1 1.2; M V & new \\
\(SU(2)_A\), \(SU(2)_B\) & --- & the two factors of the internal
symmetry of the fields that do not depend on the four extra coordinates
(zero-mode sector), in both spacetimes;
\(SO(4)=[SU(2)_A\times SU(2)_B]/Z_2\) & S1 2.1; S4 2.4; M V & new \\
\(SU(2)_X\), \(SU(2)_Y\) & --- & the two factors of the rotation group
of the four space coordinates \(x0,\dots,x3\) of the 4+4 spacetime (NOT
the internal \(SU(2)_A\times SU(2)_B\)) & DM 3.2, App. A & new \\
\(SO(3)_T\), \(SU(2)_T\) & --- & the rotations of the three additional
times \(x5,x6,x7\) of 4+4, and its covering group; the diagonal subgroup
of \(SU(2)_A\times SU(2)_B\) & S1 6.3; M V, X & new \\
\(U(1)_E\), \(P\) (generator) & --- & the second central phase symmetry
of the zero-mode sector and its traceless generator, \(P^2=1\): \(P=+1\)
on the doublets of \(SU(2)_A\), \(-1\) on those of \(SU(2)_B\) & S4 2.4;
M V & new \\
\(\Delta_{1644}^{A}\), \(\Delta_{1671}^{A}\) & \texttt{DeltaV1644}
\texttt{DeltaV1671} & the vector pair operators
\(\psi^TC^{-}\gamma^A\psi\) of each spacetime & DM 3; M IV & new \\
\(\Delta_{1671}^{AB}\) & \texttt{DeltaT1671} & the two-form pair
operators \(\psi^TC^{+}\gamma^A\gamma^B\psi\) of the 7+1 spacetime & DM
3.2 & new \\
\(\Delta^a\) & --- & order parameter: expectation value of the vector
pair operator along the extra axis \(a\in\{0,5,6,7\}\);
\(\lvert\Delta\rvert^2\equiv\sum_a\lvert\Delta^a\rvert^2\) (called \(S\)
in S1 and S4) & S1 6.1; M VIII & new \\
\(\Phi_X(\vec p)\) & --- & two-particle content
\(U(\vec p)^TXU(-\vec p)\) of the pair operator \(\psi^TX\psi\) & S1
4.1; M VI & new \\
\(\mathsf M\) & --- & the \(8\times8\) gap matrix on the internal states
of a pair & S1 6.1; M VIII & new \\
\(d_\pm\) & --- & the two fourfold gaps of a general state,
\(d_\pm^2=\lvert\Delta\rvert^2\pm\sqrt{\lvert\Delta\rvert^4-\lvert\eta_{ab}\Delta^a\Delta^b\rvert^2}\)
& S4 4.2; M VIII & new \\
\(n^a\) & --- & real unit vector: direction of the order parameter in
the space of the four extra axes & S1 6.3; M VIII & new \\
\end{longtable}

\hypertarget{the-vorton-scalar}{%
\subsection{3.6 The Vorton-Scalar}\label{the-vorton-scalar}}

\begin{longtable}[]{@{}
  >{\raggedright\arraybackslash}p{(\columnwidth - 8\tabcolsep) * \real{0.1446}}
  >{\raggedright\arraybackslash}p{(\columnwidth - 8\tabcolsep) * \real{0.1325}}
  >{\raggedright\arraybackslash}p{(\columnwidth - 8\tabcolsep) * \real{0.3976}}
  >{\raggedright\arraybackslash}p{(\columnwidth - 8\tabcolsep) * \real{0.1205}}
  >{\raggedright\arraybackslash}p{(\columnwidth - 8\tabcolsep) * \real{0.2048}}@{}}
\toprule\noalign{}
\begin{minipage}[b]{\linewidth}\raggedright
Symbol
\end{minipage} & \begin{minipage}[b]{\linewidth}\raggedright
Name in files
\end{minipage} & \begin{minipage}[b]{\linewidth}\raggedright
Meaning
\end{minipage} & \begin{minipage}[b]{\linewidth}\raggedright
Defined in
\end{minipage} & \begin{minipage}[b]{\linewidth}\raggedright
Origin
\end{minipage} \\
\midrule\noalign{}
\endhead
\bottomrule\noalign{}
\endlastfoot
\(\Phi_{VS}\) & \texttt{PhiVS} & the Vorton-Scalar: a complex scalar
field & S1 1.2; M V & new (the name is that of the author's project
NCSS, where the scalar is real) \\
\(v_{VS}\) & \texttt{vVS} & vacuum value:
\(\langle\Phi_{VS}\rangle=v_{VS}/\sqrt2\) & S1 1.2; M V & new \\
\(\lambda_{VS}\) & \texttt{lambdaVS} \texttt{lamVS} & self-coupling of
the Vorton-Scalar & S1 1.2; M V & new \\
\(y_{VS}\) & \texttt{yVS} & Yukawa coupling of the Vorton-Scalar to the
fermion field & S1 1.2; M V & new \\
\(s_{VS}\) & \texttt{sVS} & the radial mode of the Vorton-Scalar (real
field) & S1 1.4; M V & new \\
\(m_s\) & \texttt{msVS} & mass of the radial mode,
\(\sqrt{2\lambda_{VS}}\,v_{VS}\) & S1 1.4; M V & new \\
\(\vartheta_{VS}\) & \texttt{thetaVS} \texttt{thVS} & the phase mode of
the Vorton-Scalar (real field, massless) & S1 1.4; M V & new \\
\(g_{VS}\) & \texttt{gVS} & coupling of both modes to the fermions,
\(m/v_{VS}=y_{VS}/\sqrt2\) & S1 1.4; M V & new \\
\(\alpha_{VS}\) & \texttt{alphaVS} \texttt{aVS} & \(g_{VS}^2/4\pi\) & S1
1.4; M V & new \\
\(D_{VS}\) & \texttt{DVS} & \(\alpha_{VS}\,m/m_s\): two fermions bind if
and only if \(D_{VS}>D_c\) & S1 5.3; M VII & new \\
\(J_{VS}^A\) & \texttt{JVS} & the neutral current of the Vorton-Scalar
(bilinear in \(\Phi_{VS}\)) & S1 1.5; M V & new \\
\(J_5^A\) & \texttt{J5} & the axial current \(\bar\psi\gamma^A\chi\psi\)
of the fermion field & S1 1.5; M V & new \\
\(G_{VS}\) & \texttt{GVS} & constant of the Fermi type of the
neutral-current interaction, \(G_{VS}/\sqrt2=1/(2v_{VS}^2)\) & S1 1.5; M
V & new \\
\(\mathcal L^{VS}_{1644}\), \(\mathcal L^{VS}_{1671}\) & \texttt{LVS} &
the Lagrangians that define the Vorton-Scalar in the two spacetimes & S1
1.2, 3.2; M V & new \\
\(C\) (coupling) & --- & \(g_{VS}^2m^2/m_s^2=g_{VS}^4/2\lambda_{VS}\),
the one number on which the fermion mass in matter depends & S4 5.1; M
IX & new \\
\end{longtable}

\hypertarget{pairing-the-gap-equation-and-the-effective-action-of-cornwall-jackiw-and-tomboulis}{%
\subsection{3.7 Pairing, the gap equation and the effective action of
Cornwall, Jackiw and
Tomboulis}\label{pairing-the-gap-equation-and-the-effective-action-of-cornwall-jackiw-and-tomboulis}}

\begin{longtable}[]{@{}
  >{\raggedright\arraybackslash}p{(\columnwidth - 8\tabcolsep) * \real{0.1446}}
  >{\raggedright\arraybackslash}p{(\columnwidth - 8\tabcolsep) * \real{0.1325}}
  >{\raggedright\arraybackslash}p{(\columnwidth - 8\tabcolsep) * \real{0.3976}}
  >{\raggedright\arraybackslash}p{(\columnwidth - 8\tabcolsep) * \real{0.1205}}
  >{\raggedright\arraybackslash}p{(\columnwidth - 8\tabcolsep) * \real{0.2048}}@{}}
\toprule\noalign{}
\begin{minipage}[b]{\linewidth}\raggedright
Symbol
\end{minipage} & \begin{minipage}[b]{\linewidth}\raggedright
Name in files
\end{minipage} & \begin{minipage}[b]{\linewidth}\raggedright
Meaning
\end{minipage} & \begin{minipage}[b]{\linewidth}\raggedright
Defined in
\end{minipage} & \begin{minipage}[b]{\linewidth}\raggedright
Origin
\end{minipage} \\
\midrule\noalign{}
\endhead
\bottomrule\noalign{}
\endlastfoot
\(\nu\) & --- & the number of internal states of a particle of given
momentum, \(\nu=8\) & S1 5; M VII & new \\
\(F_S\), \(F_P\) & \texttt{F\_scalar} \texttt{F\_pseudoscalar} &
exchange amplitudes of the radial and of the phase mode in a pair
channel & S1 4.1; M VI & new \\
\(V_{\rm NC}\), \(V_s(r)\) & --- & contact interaction of the neutral
current in the hidden-axis channel; Yukawa potential of the radial mode
& S1 4.2; M VI & new \\
\(\sigma_X\) & --- & the sign \(+1\) (vector pairs along the extra axes)
or \(-1\) (their two-form pairs) of the phase-mode term of the kernel &
M VII & new \\
\(W(k,k')\) & --- & \(s\)-wave average of the exchange attraction:
kernel of the gap equation & S1 5.1; M VII & new \\
\(\Delta(k)\), \(\Delta_F\) & --- & gap function; its value at the Fermi
momentum & S1 5.1; M VII & new \\
\(\xi_k\), \(E_k\), \(\kappa(k)\) & --- & \(\sqrt{m^2+k^2}-\mu\);
quasi-particle energy \(\sqrt{\xi_k^2+\Delta(k)^2}\); pair correlation
\(\langle b(-\vec k)b(\vec k)\rangle\) & S4 3.1; M VII & new \\
\(k_F\), \(E_F\), \(\mu\) & --- & Fermi momentum,
\(n=\nu k_F^3/6\pi^2\); Fermi kinetic energy \(\sqrt{m^2+k_F^2}-m\);
chemical potential & S1 5; M VII & new \\
\(N(0)\) & --- & density of states at the Fermi surface per internal
state, \(k_F\sqrt{m^2+k_F^2}/2\pi^2\) & S1 5.2; M VII & new \\
\(\lambda\), \(\lambda_S\), \(\lambda_P\) & --- & dimensionless pairing
coupling \(N(0)W(k_F,k_F)\) and its parts from the radial and the phase
mode & S1 5.2; M VII & new \\
\(T_c\) & --- & critical temperature of the paired state & S1 5.2; M VII
& new \\
\(D_c\) & --- & critical value of \(D_{VS}\) for a bound pair in the
Yukawa potential & S1 5.3; M VII & new \\
\(a_{\rm sc}\), \(a_{BB}\) & --- & scattering length of two fermions
(written \(a\) in S1, S3, S4); scattering length of two bound pairs & S1
5.3; M VII, VIII & new \\
\(\Gamma[\phi,G]\), \(\Gamma_2\) & --- & the effective action for
composite operators of Cornwall, Jackiw and Tomboulis; its
two-particle-irreducible part & S4 1; M VII & new (method of the cited
paper) \\
\(\mathcal S\), \(D_s\), \(D_\vartheta\) & --- & propagator of the
doubled fermion field; propagators of the radial and of the phase mode &
S4 2.1; M VII & new \\
\(V[\Delta]\) & --- & the CJT effective potential of the paired state,
measured from the unpaired state & S4 3.1; M VII & new \\
\(\eta(T)\) & --- & largest eigenvalue of the linearized gap equation,
\(\eta(T_c)=1\) & S4 3.4; M VII & new \\
\(A_2(T)\), \(A_4\) & --- & coefficients of the expansion of \(V\) along
a gap function of fixed shape (written \(\alpha\), \(\beta\) in S4) & S4
3.4; M VII & new \\
\(a(T)\), \(b\), \(K_{\rm GL}\) & --- & Ginzburg-Landau coefficients of
the weak-coupling free energy (\(K_{\rm GL}\) is written \(K\) in S1 and
S4) & S1 6.2; M VIII & new \\
\(c_s\) & --- & speed of first sound of the superfluid,
\(c_s^2=dP/d\rho\) & S1 6.4; M VIII & new \\
\end{longtable}

\hypertarget{matter-temperature-and-vacuum}{%
\subsection{3.8 Matter, temperature and
vacuum}\label{matter-temperature-and-vacuum}}

\begin{longtable}[]{@{}
  >{\raggedright\arraybackslash}p{(\columnwidth - 8\tabcolsep) * \real{0.1446}}
  >{\raggedright\arraybackslash}p{(\columnwidth - 8\tabcolsep) * \real{0.1325}}
  >{\raggedright\arraybackslash}p{(\columnwidth - 8\tabcolsep) * \real{0.3976}}
  >{\raggedright\arraybackslash}p{(\columnwidth - 8\tabcolsep) * \real{0.1205}}
  >{\raggedright\arraybackslash}p{(\columnwidth - 8\tabcolsep) * \real{0.2048}}@{}}
\toprule\noalign{}
\begin{minipage}[b]{\linewidth}\raggedright
Symbol
\end{minipage} & \begin{minipage}[b]{\linewidth}\raggedright
Name in files
\end{minipage} & \begin{minipage}[b]{\linewidth}\raggedright
Meaning
\end{minipage} & \begin{minipage}[b]{\linewidth}\raggedright
Defined in
\end{minipage} & \begin{minipage}[b]{\linewidth}\raggedright
Origin
\end{minipage} \\
\midrule\noalign{}
\endhead
\bottomrule\noalign{}
\endlastfoot
\(\phi\) & --- & expectation value of the Vorton-Scalar,
\(\langle\Phi_{VS}\rangle=\phi/\sqrt2\); the fermion mass is
\(g_{VS}\phi\) & S4 2.1; M VII & new \\
\(m^*\), \(x\) & --- & fermion mass in the medium, and \(x=m^*/m\) (in
the manuscripts Sec. XII A, \(x\) denotes \(k_F/m\), stated there) & S4
5.1; M IX & new \\
\(\varepsilon(k_F,x)\), \(V_{\rm sea}(x)\) & --- & energy density of the
unpaired state in units of \(m^4\); renormalized energy of the filled
negative-energy states & S4 5.1; M IX & new \\
\(B\) & --- & energy density of the state without expectation value of
the Vorton-Scalar (bag constant) & S4 5.2; M IX & new \\
\(E_x\) & --- & exchange (Fock) energy of the two modes & S4 5.4; M IX &
new \\
\(M_s\), \(M_\vartheta\), \(a_s\), \(a_\vartheta\) & --- & thermal
masses of the two boson modes in the Hartree approximation, and the
curvatures \(U_{\rm eff}''\), \(U_{\rm eff}'/\phi\) & S4 6; M IX &
new \\
\(P(M)\), \(Q(M)\) & --- & thermal part of the boson loop and of
\(\tfrac12\mathrm{Tr}\ln\) & S4 6; M IX & new \\
\(T_{\rm high}\), \(T_N\), \(T_{1/2}\), \(T_0\) & --- & two closed
estimates of the temperature at which the fermion mass vanishes;
temperatures at which the expectation value has fallen to half and to
zero & S4 6; M IX & new \\
\(\sigma(p)\) & --- & self-energy that a pair condensate in the vacuum
gives to the fermions: mass splitting \(m\pm\sigma\) & S4 7.2; M IX &
new \\
\(\alpha_c\), \(\Lambda\) (cutoff) & --- & critical value of
\(\alpha_{VS}\) for a pair condensate in the vacuum; momentum cutoff of
the ladder equation & S4 7.2; M IX & new \\
\end{longtable}

\hypertarget{cosmology}{%
\subsection{3.9 Cosmology}\label{cosmology}}

\begin{longtable}[]{@{}
  >{\raggedright\arraybackslash}p{(\columnwidth - 8\tabcolsep) * \real{0.1446}}
  >{\raggedright\arraybackslash}p{(\columnwidth - 8\tabcolsep) * \real{0.1325}}
  >{\raggedright\arraybackslash}p{(\columnwidth - 8\tabcolsep) * \real{0.3976}}
  >{\raggedright\arraybackslash}p{(\columnwidth - 8\tabcolsep) * \real{0.1205}}
  >{\raggedright\arraybackslash}p{(\columnwidth - 8\tabcolsep) * \real{0.2048}}@{}}
\toprule\noalign{}
\begin{minipage}[b]{\linewidth}\raggedright
Symbol
\end{minipage} & \begin{minipage}[b]{\linewidth}\raggedright
Name in files
\end{minipage} & \begin{minipage}[b]{\linewidth}\raggedright
Meaning
\end{minipage} & \begin{minipage}[b]{\linewidth}\raggedright
Defined in
\end{minipage} & \begin{minipage}[b]{\linewidth}\raggedright
Origin
\end{minipage} \\
\midrule\noalign{}
\endhead
\bottomrule\noalign{}
\endlastfoot
\(N\) & --- & density of the dark radiation today in units of the
density of one massless neutrino species & S3 3.2; M XI & new \\
\(z_t\), \(T_d\) & --- & transition redshift of the stepped dark fluid,
\(m=T_{d0}(1+z_t)\); temperature of the dark sector & S3 3.2; M XI, D &
new \\
\(r_{\rm step}\) & --- & factor by which the density of the dark
radiation rises when the fermions annihilate, \(15^{1/3}\) for this
model & S3 1; M XI & new \\
\(\omega_b\), \(\omega_c\), \(h\), \(H_0\), \(A_s\), \(n_s\), \(\tau\) &
--- & the six parameters of the standard cosmological model: densities
of baryons and of cold dark matter, Hubble constant \(H_0=100\,h\) km
s\(^{-1}\) Mpc\(^{-1}\), amplitude and tilt of the primordial spectrum,
optical depth & S3 2.1; M XI & standard notation \\
\(\Delta\chi^2\) & --- & worsening of the best fit to the mock data
relative to the standard model & S3 3.3; M XI & new \\
\(\theta_s\), \(r_s\), \(D_A\) & --- & angular size of the sound
horizon; sound horizon; distance to last scattering & S3 3.1; M XI &
standard notation \\
\(c_{s0}\) & --- & sound speed of superfluid dark matter today at the
mean cosmic density & S3 4.1; M XII & new \\
\(\rho_0\), \(R_{TF}\) & --- & constant of the equation of state of a
condensate of bound pairs; Thomas-Fermi radius of a core supported by
its pressure & S3 4.1; M XII & new \\
\(m_{\rm deg}\), \(m_{B,\max}\) & --- & largest fermion mass for
degeneracy, and largest boson mass for condensation, at a given density
and velocity dispersion & S3 4.3; M XII & new \\
\end{longtable}

\begin{center}\rule{0.5\linewidth}{0.5pt}\end{center}

\hypertarget{letters-that-have-more-than-one-meaning-in-the-documents-of-the-project}{%
\section{4. Letters that have more than one meaning in the documents of
the
project}\label{letters-that-have-more-than-one-meaning-in-the-documents-of-the-project}}

The tagged symbols of Sec. 3 have one meaning each. The plain letters
below are used for more than one quantity. In every case the meanings
are told apart by a subscript, a superscript, an argument, a different
typeface, or by the section in which the letter occurs; the last column
says which. The table covers the two manuscripts (M, with the Roman
number of the section), the four supplements and ``dirac matrices''
(DM).

The table was written by hand and then checked against a list made by a
program: \texttt{data/symbol\_concordance.csv} (script
\texttt{code/scripts/symbol\_concordance.py}) gives, for every Latin and
Greek letter and every document, each form in which the letter occurs in
a formula, with the number of occurrences by section. A reader who wants
to see every use of a letter finds it there.

\begin{longtable}[]{@{}
  >{\raggedright\arraybackslash}p{(\columnwidth - 4\tabcolsep) * \real{0.0706}}
  >{\raggedright\arraybackslash}p{(\columnwidth - 4\tabcolsep) * \real{0.6000}}
  >{\raggedright\arraybackslash}p{(\columnwidth - 4\tabcolsep) * \real{0.3294}}@{}}
\toprule\noalign{}
\begin{minipage}[b]{\linewidth}\raggedright
Letter
\end{minipage} & \begin{minipage}[b]{\linewidth}\raggedright
Meanings
\end{minipage} & \begin{minipage}[b]{\linewidth}\raggedright
How they are told apart
\end{minipage} \\
\midrule\noalign{}
\endhead
\bottomrule\noalign{}
\endlastfoot
\(a\), \(b\) & index of the four extra axes, values 0, 5, 6, 7 (S1 to
S4; M IV to XIII); index of the sixteen components of the field,
\(\psi_a\) (M III, IV); in DM Sec. 1.8.4, index of the four space-like
axes of 4+4; scale factor \(a\) of the universe (M XII); Ginzburg-Landau
coefficients \(a(T)\), \(b\); curvatures \(a_s\), \(a_\vartheta\) of the
potential; scattering length (S1, S3, S4); annihilation operators
\(b(\vec p)\) & an index is a subscript or superscript, and its range is
stated where it is used; \(a(T)\) carries its argument, \(a_s\),
\(a_\vartheta\) their subscript; the manuscripts write the scattering
length \(a_{\rm sc}\); an operator carries its argument \\
\(A\), \(B\) & indices with the values 0 to 7; the coefficients
\(A_2(T)\), \(A_4\); the amplitude \(A_s\) of the primordial spectrum;
the bag constant \(B\) (S4; M IX); the Killing form \(B(X,Y)\) (DM; M
II); subscripts of the groups \(SU(2)_A\), \(SU(2)_B\) and of quantities
of a bound pair (\(m_B\), \(n_B\), \(a_{BB}\)) & subscript; argument;
position \\
\(C\) & index; the matrices \(C^{-}\), \(C^{+}\); the coupling
\(C=g_{VS}^2m^2/m_s^2\); heat capacities \(C_s\), \(C_n\); angular power
spectra \(C_\ell\); the labels C1 to C7 of the conditions on a Higgs
field & superscript \(\mp\) and tag; subscript; a label is followed by a
digit \\
\(D\) & index; \(D_{VS}\) and its critical value \(D_c\); inverse
propagator \(D^{-1}(\phi)\) and the boson propagators \(D_s\),
\(D_\vartheta\); \(D_\ell\equiv\ell(\ell+1)C_\ell/2\pi\); distance
\(D_A\) & subscript \\
\(d\) & the two gaps \(d_\pm\) of a general state; annihilation
operators \(d_s(\vec p)\) of antiparticles (M III); differential &
subscript and argument \\
\(E\) & energies \(E\), \(E_k\), \(E_F\), \(E_b\), \(E_x\); the chiral
frames \(E_{1671}^{L,R}\); subscript of the group \(U(1)_E\) & tag and
superscript; position \\
\(G\) & full propagator \(G\) of the effective action (M VII, Appendix
B; S4 Sec. 1); the constant of gravitation (M XII; S3 Sec. 4.1);
\(G_{VS}\); Fermi constant \(G_F\) & subscript; section \\
\(\Gamma\) & effective action \(\Gamma[\phi,G]\) and its part
\(\Gamma_2\) (M VII; S4); a vertex matrix \(\Gamma\) in the exchange
amplitude (M VI; S1 Sec. 4) & argument or subscript; section \\
\(h\) & index of the matrices \(s_h\), \(t_h\); Hubble constant in units
of 100 km s\(^{-1}\) Mpc\(^{-1}\) (M XI, XII; S3); Hamiltonian matrix
\(\mathsf h\) & index position; the matrix is set in sans-serif type and
carries a tag where it is defined \\
\(K\) & \(K_{1644}\), generator of the evolution in \(x4\) (DM; M III of
4+4); the source \(K(x,y)\) of the effective action; gradient
coefficient of the Ginzburg-Landau energy (S1, S4) & tag; argument; the
manuscripts write the gradient coefficient \(K_{\rm GL}\) \\
\(M\) & mass of an exchanged boson (M VI, IX; S4); thermal masses
\(M_s\), \(M_\vartheta\); gap matrix \(\mathsf M\); the matrix
\(M_{1671}\) of a worked example (DM; M III of 7+1); \(M_W\), \(M_Z\);
solar mass \(M_\odot\) & subscript; typeface; tag \\
\(N\) & density of dark radiation (M XI; S3); density of states
\(N(0)\); norms \(N_X\), \(N_a\) of pair amplitudes; \(N_{\rm eff}\);
subscript of the temperature \(T_N\) & argument; subscript \\
\(n\) & number density; direction \(n^a\) of the order parameter;
winding number of a vortex (M V; S1); tilt \(n_s\); size of the matrices
\(1_n\) and of the spectra in the mock likelihood & index; subscript;
section \\
\(P\) & chiral projectors \(P^{L}\), \(P^{R}\); the central generator
\(P\) of \(U(1)_E\) (M V; S4); pressure \(P\), \(P(\rho)\); thermal
boson loops \(P(M)\), \(P_s\), \(P_\vartheta\); subscript of the
amplitude \(F_P\) and of the coupling \(\lambda_P\) of the phase mode &
superscript and tag; argument; subscript; section \\
\(Q\) & thermal part \(Q(M)\) of \(\tfrac12{\rm Tr}\ln\) (M IX; S4);
electric charge (M X; S2) & argument; section \\
\(r\), \(s\) & labels of the eight internal states of a particle,
\(u_s\), \(b_r\) (M III, VI); rank \(r\) of a product of Dirac matrices
(DM; M II, IV); radius \(r\); the matrices \(s_h\); the radial mode
\(s_{VS}\) and the subscript \(s\) of its quantities (\(m_s\), \(D_s\),
\(M_s\), \(a_s\), \(P_s\)); scaling factor \(s\) of a trial gap function
(M VII; S4); subscript of the superfluid state (\(S_s\), \(C_s\));
standard cosmological subscripts (\(c_s\), \(r_s\), \(\theta_s\),
\(n_s\), \(A_s\)); \(r_{\rm step}\), \(r_{\rm drag}\) & index position;
suffix VS; the quantity that carries the subscript \\
\(S\) & element \(S\) of the spin group (DM; M II to IV); classical
action \(S[\Phi]\); entropy \(S\), \(S_s\), \(S_n\); propagator
\(\mathcal S\) of the doubled fermion field; subscript of \(F_S\),
\(\lambda_S\) (radial mode); in S1 and S4,
\(S=\sum_a\lvert\Delta^a\rvert^2\) & argument; subscript; typeface; the
manuscripts write \(\lvert\Delta\rvert^2\) for the last \\
\(\Sigma\) & the scaled commutators \(\Sigma^{AB}\); self-energy
\(\Sigma\), \(\Sigma[\phi,G]\) of the effective action (M VII, Appendix
B; S4) & indices \\
\(T\) & temperature, with \(T_c\), \(T_{1/2}\), \(T_0\),
\(T_{\rm high}\), \(T_N\), \(T_d\); transposition (superscript);
generators \(T^3_A\), \(T^3_B\) of \(SU(2)_A\), \(SU(2)_B\) (M X; S2);
subscript of the group \(SO(3)_T\) of the additional times & position;
indices \\
\(t\) & index of the three additional coordinates, values 5, 6, 7; the
matrices \(t_h\); subscript of the transition redshift \(z_t\) & index
position \\
\(U\) & matrix \(U(\vec p)\) of positive-energy spinors;
\(U_{1671}^{c}\); the effective potential \(U_{\rm eff}\); the groups
\(U(1)\) & argument; tag; subscript \\
\(u\), \(v\) & the real spinors \(u_s\), \(v_s\) of the plane waves (M
III); reduced radial wave function \(u\) of the bound pair (M VII; S1);
the vacuum value \(v_{VS}\) and the electroweak value \(v_{\rm EW}\);
Fermi velocity \(v_F\); a vector \(v\) in the definition of the Clifford
algebra (M II) & subscript; section \\
\(V\) & effective potential \(V\), \(V[\Delta]\), \(V(\phi,G)\); the
interactions \(V_{\rm NC}\), \(V_s(r)\), \(V_{\rm BCS}\);
\(V_{\rm sea}\); the matrices \(V_1\), \(V_2\) of the exchange amplitude
(M VI); volume \(V\) of the quantization box (M III) and \(V_4\) of the
extra dimensions & subscript; argument; section \\
\(W\) & \(W_{1644}\), the grading; kernel \(W(k,k')\) of the gap
equation; generating functional \(W[J,K]\); subscript of the weak mixing
angle \(\theta_W\) and of \(M_W\) & tag; argument; position \\
\(x\), \(y\) & the coordinates \(x0,\dots,x7\) and points \(x\), \(y\)
of spacetime; the ratio \(x=m^*/m\) (M IX; S4); in M XII A, \(x=k_F/m\);
the Yukawa coupling \(y_{VS}\) & the coordinates carry a digit; the two
ratios are defined where they are used; suffix VS \\
\(\alpha\) & the matrices \(\alpha^j\); \(\alpha_{VS}\) and its critical
value \(\alpha_c\); in S4, a coefficient \(\alpha(T)\); in the name
``Lyman-\(\alpha\)'' & index; subscript; the manuscripts write
\(A_2(T)\) for the coefficient \\
\(\beta\) & the matrix \(\beta\); in S4, a coefficient & the manuscripts
write \(A_4\) for the coefficient \\
\(\Delta\) & order parameter \(\Delta^a\) and gap \(\Delta(k)\),
\(\Delta_F\); the pair operators \(\Delta_{1644}^A\),
\(\Delta_{1671}^A\); differences \(\Delta\chi^2\),
\(\Delta N_{\rm eff}\) & a difference is followed by the quantity that
differs \\
\(\epsilon\), \(\varepsilon\) & parameter of a symmetry transformation;
permutation symbol \(\epsilon_{hpq4}\); energy density
\(\varepsilon(k_F,x)\); dimensionless binding energy of the pair (M VII
E; S1) & indices; argument; section \\
\(\eta\) & the metrics \(\eta_{1644}\), \(\eta_{1671}\), \(\eta^{AB}\),
\(\eta_{ab}\); eigenvalue \(\eta(T)\) of the linearized gap equation &
indices or tag; argument \\
\(\theta\), \(\vartheta\) & rotation parameters \(\theta_{AB}\),
\(\theta\); angle between two momenta (M VII) and angle of a direction
of the order parameter (M VIII); angular size \(\theta_s\) of the sound
horizon; weak mixing angle \(\theta_W\); the phase mode
\(\vartheta_{VS}\) and the subscript \(\vartheta\) of its quantities &
subscript; suffix VS; section \\
\(\kappa\) & pair correlation \(\kappa(k)\); frequencies \(\kappa_t\)
along the additional times of 4+4 (M III of 4+4; DM); the bound
\(\kappa=1\) cm\(^2\)/g on the self-interaction (M XII; S3) & argument;
index; unit \\
\(\Lambda\) & Lorentz matrix \(\Lambda^A{}_B\); momentum cutoff; the
cosmological constant in the name ΛCDM & indices; context \\
\(\lambda\) & \(\lambda_{VS}\); pairing couplings \(\lambda\),
\(\lambda_S\), \(\lambda_P\); the coupling of the zero-dimensional test
model (M Appendix B; S4 Sec. 1.7); thermal wavelength (M XII) & suffix
VS; subscript; section \\
\(\mu\), \(\nu\) & chemical potential \(\mu\); the number \(\nu=8\) of
internal states; in the unit \(\mu\)K & context \\
\(\rho\) & mass density, with \(\rho_0\), \(\rho_{\rm deg}\); the
variable \(m_sr\) of the bound-pair equation (M VII E; S1 Sec. 5.3); the
ratio \(M_W^2/M_Z^2\cos^2\theta_W\) (M X; S2); an exponent in the
solution of the ladder equation (M IX) & defined where used; each of the
last three occurs in one subsection \\
\(\sigma\) & the matrix \(\sigma_{88}\); the sign \(\sigma_X\);
self-energy \(\sigma\), \(\sigma(p)\) (M IX; S4); velocity dispersion (M
XII; S3); statistical errors \(\sigma(N)\), \(\sigma_\tau\);
\(\sigma_8\) & subscript; argument; section \\
\(\tau\) & the matrices \(\tau_{1644}^A\), \(\tau_{1671}^A\); optical
depth of reionization, with its prior value \(\tau_0\) & tag and index;
the optical depth occurs in the cosmological sections only \\
\(\phi\), \(\varphi\) & expectation value \(\phi\) of the Vorton-Scalar;
chiral parts \(\phi_{1671}^{L,R}\); shape \(\varphi(k)\) of the gap
function; phase \(\varphi\) of the order parameter; variable of the
zero-dimensional test model & tag and superscript; argument; section \\
\(\Phi\) & \(\Phi_{VS}\); two-particle content \(\Phi_X(\vec p)\),
\(\Phi_a(\vec p)\); the generic field \(\Phi\) of the general method (M
VII A; S4 Sec. 1) & suffix VS; subscript \\
\(\chi\) & chirality matrix; the statistic \(\chi^2\) & exponent 2 \\
\(\omega\) & \(\omega_{1644}\), \(\omega_{1671}\); the densities
\(\omega_b\), \(\omega_c\) (M XI; S3); a cutoff frequency \(\omega_c\)
in the textbook formula for the gap (M VII; S1 Sec. 5.2) & tag;
subscript and section \\
\end{longtable}

Seven letters were changed in the documents while this provenance was
written, because two meanings occurred close to each other: in the
manuscripts the sum \(\sum_a\lvert\Delta^a\rvert^2\) is written
\(\lvert\Delta\rvert^2\) (the supplements call it \(S\), which the
manuscripts need for the propagator \(\mathcal S\)); the scattering
length is written \(a_{\rm sc}\); the expansion coefficients of the
effective potential are written \(A_2\), \(A_4\) (the supplement S4
calls them \(\alpha\), \(\beta\), the letters of the Hamiltonian
matrices); the gradient coefficient is written \(K_{\rm GL}\); the index
of the four space axes in the evolution formula of 4+4 is \(j\), not
\(a\); the matrix of the mock data in the likelihood is written
\(\hat C_\ell\), not \(D_\ell\); and, in the manuscripts and in S1, the
generator of an internal symmetry is written \(\mathsf g\), not \(a\).

\begin{center}\rule{0.5\linewidth}{0.5pt}\end{center}

\hypertarget{the-original-files-one-by-one}{%
\section{5. The original files, one by
one}\label{the-original-files-one-by-one}}

For each file: what it is, which conventions it uses, which of its
symbols collide with symbols of other files, and what the project does.
The provenances P2 (papers) and P3 (notebooks) state how each file was
used. Every statement about a notebook in this section and the next was
read off the notebook itself; for the symbols, the defining statements
were collected by the program of Appendix A.

\hypertarget{the-notebooks}{%
\subsection{5.1 The notebooks}\label{the-notebooks}}

\begin{longtable}[]{@{}
  >{\raggedright\arraybackslash}p{(\columnwidth - 6\tabcolsep) * \real{0.0682}}
  >{\raggedright\arraybackslash}p{(\columnwidth - 6\tabcolsep) * \real{0.7045}}
  >{\raggedleft\arraybackslash}p{(\columnwidth - 6\tabcolsep) * \real{0.1136}}
  >{\raggedleft\arraybackslash}p{(\columnwidth - 6\tabcolsep) * \real{0.1136}}@{}}
\toprule\noalign{}
\begin{minipage}[b]{\linewidth}\raggedright
Id
\end{minipage} & \begin{minipage}[b]{\linewidth}\raggedright
Notebook
\end{minipage} & \begin{minipage}[b]{\linewidth}\raggedleft
Symbols defined
\end{minipage} & \begin{minipage}[b]{\linewidth}\raggedleft
of these, also defined in another notebook
\end{minipage} \\
\midrule\noalign{}
\endhead
\bottomrule\noalign{}
\endlastfoot
N01 & \texttt{4+4.nb} & 144 & 142 \\
N02 & \texttt{7+1.nb} & 1429 & 86 \\
N03 & \texttt{Bessel\ product.nb} & 6 & 0 \\
N04 & \texttt{Dirac\_Einstein-Lovelock\_4+4\_Nash.nb} & 80 & 79 \\
N05 & \texttt{Dirac\_Nash\_Identity.nb} & 82 & 41 \\
N06 & \texttt{Einstein-Lovelock-4+4-Nash.nb} & 69 & 69 \\
N07 & \texttt{Einstein-Lovelock-7+1-Nash.nb} & 68 & 68 \\
N08 & \texttt{Einstein-Lovelock-Eternal.nb} & 67 & 59 \\
N09 & \texttt{Generalized\ \_Kronecker\_Delta\_4+4.nb} & 68 & 55 \\
N10 &
\texttt{Pair\_Creation\_of\_Universes\_WaveFunctionOfUniverse-4+4-Einstein-Lovelock-Nash.nb}
& 368 & 278 \\
N11 & \texttt{WaveFunctionUniverse-4+4-Einstein-Lovelock-Nash.nb} & 193
& 174 \\
N12 & \texttt{WaveFunctionUniverse-4+4-Einstein-Lovelock.nb} & 247 &
200 \\
N13 & \texttt{det-g.nb} & 6 & 5 \\
N14 & \texttt{field\_equations.nb} & 60 & 60 \\
\end{longtable}

The numbers in the table come from the scan described in Appendix A.

\textbf{N01, \texttt{4+4.nb}} (title: ``eight, real-valued, 16 by 16,
Dirac Matrices''). The source of the 4+4 matrices. \emph{Conventions:} a
text at its head gives the coordinates
\(\{x0;\,x1,x2,x3;\,x4;\,x5,x6,x7\}\) as hidden space, 3-space, time and
``superluminal deflating time''; index \(A=0,\dots,7\); metric
\texttt{η4488} \(=\) diag\((1,1,1,1,-1,-1,-1,-1)\), space-like \(=+1\).
With the package Notation the notebook turns typeset expressions into
names: \texttt{τ} with an overbar, \texttt{T16\^{}A},
\texttt{F\_A\^{}a}. \emph{Symbols taken over:} \texttt{τ{[}A{]}} and
\texttt{τ} with an overbar \(\to\tau_{1644}^A\), \(\bar\tau_{1644}^A\);
\texttt{σ} \(=\) \texttt{σ88} \(\to\sigma_{88}\);
\texttt{T16\^{}A{[}A{]}} \(\to\gamma_{1644}^A\);
\texttt{T16\^{}A{[}8{]}} \(\to\omega_{1644}\); \texttt{P\_L},
\texttt{P\_R} \(\to P_{1644}^{L,R}\); \texttt{σ16} \(\to C_{1644}^{-}\);
\texttt{η4488} \(\to\eta_{1644}\); \texttt{s4by4{[}h{]}},
\texttt{t4by4{[}h{]}} \(\to s_h\), \(t_h\). \emph{Collisions:}
\texttt{γ} is a list of sixteen \(4\times4\) matrices (four matrices in
N05, a tensor index in N09, the \(16\times16\) Dirac matrices in the
project); \texttt{σ}, \texttt{τ}, \texttt{ζ} are matrices here and
tensor indices in N09; \texttt{h} is the number 8 that selects a unit
spinor (\texttt{h\ =\ hUSE}), \texttt{Ω} is the matrix
\texttt{σ.τ{[}7{]}}, \texttt{J} and \texttt{B} are auxiliary. Inside
\texttt{Block{[}...{]}} the notebook also uses the letters \texttt{m},
\texttt{ψ}, \texttt{t}, \texttt{b}, \texttt{r}, \texttt{s} as temporary
variables (\texttt{m}, for example, for a \(16\times16\) array of
unknowns); these are local and are not counted as symbols of the
notebook. \emph{Remedy:} only the tagged symbols are used; the auxiliary
variables are not used at all.

\textbf{N02, \texttt{7+1.nb}} (title: ``7+1: the octonions and the
Clifford algebra of 7+1 spacetime''; 27 sections). The source of the 7+1
matrices. The notebook says that it was generated by a script for the
author's 7+1 project and evaluated in Mathematica; it reproduces the
objects of N01 under plain names and builds the 7+1 objects from them.
\emph{Conventions:} the same eight coordinate names as N01, but with
\(x0\) as the time (its Sec. 1.1 says so in a table); two sign
conventions side by side (its Sec. 1.2), \texttt{η71} \(=\)
diag\((-1,1,\dots,1)\), which is the convention of N01, and \texttt{η17}
\(=-\)\texttt{η71}. \emph{Symbols taken over:} \texttt{Γ17up{[}A{]}}
\(\to\gamma_{1671}^{\mathrm{swap}_{04}(A)}\); \texttt{Γ17{[}8{]}},
\texttt{ω17} \(\to\omega_{1671}\); \texttt{χ17} \(\to\chi_{1671}\);
\texttt{PL17}, \texttt{PR17} \(\to P_{1671}^{L,R}\); \texttt{τO{[}A{]}}
\(\to\tau_{1671}^{\mathrm{swap}_{04}(A)}\); \texttt{βD},
\texttt{αD{[}i{]}} \(\to\beta_{1671}\), \(\alpha_{1671}^{j}\);
\texttt{η17} \(\to\eta_{1671}\) after the exchange of the index values 0
and 4; \texttt{S44} \(\to C_{1644}^{+}\); \texttt{dirac16complex}
\(\to\psi_{1671}\) (dirac1671); \texttt{EL}, \texttt{ER}
\(\to E_{1671}^{L,R}\). \emph{Collisions inside the file:} \texttt{Γ17}
(real, its Sec. 4.2) against \texttt{Γ71} (complex, Sec. 4.3) and
\texttt{Γ71R32} (real \(32\times32\)); \texttt{η17} against
\texttt{η71}; \texttt{χ17} against \texttt{χ71}; \texttt{PL17},
\texttt{PR17} against \texttt{PL71}, \texttt{PR71} and \texttt{PL44},
\texttt{PR44}; \texttt{τ} (4+4) against \texttt{τO} (octonionic).
Several letters change their meaning from section to section: \texttt{H}
is the Hamiltonian of the field (Sec. 7.7), a flavour Hamiltonian (Sec.
19) and a constant declared for the tensor package (Sec. 12.2);
\texttt{M} is an intertwining matrix (Sec. 4.4), the \(8\times8\) matrix
\(m+p_i\,\)\texttt{τO{[}i{]}} (Sec. 7.4) and a Kohn--Sham mass (Sec.
9.3); \texttt{S} is a unitary matrix (Sec. 4.3), a scalar potential
(Sec. 8.4) and the action (Secs. 12, 17); \texttt{κ} is a chiral density
(Sec. 9.1) and the constant \(8\pi G_8\) (Sec. 11); \texttt{h} is the
Hamiltonian matrix (Sec. 7.3) and an induced metric (Secs. 20.3, 25.6);
\texttt{m} is the mass of the field and, for the tensor package, an
abstract index. \emph{Collisions with other files:} \texttt{x0} is the
time here and a hidden space coordinate in N01 and N04; the dictionary
of its Sec. 1.3 renames objects of N01 (Sec. 6.2 below); \texttt{sAB} is
a sign here and a generator in N10 to N12. \emph{Remedy:} the project
uses the real set only, under the tag 1671, with the time at index 4;
the relabelling \(\mathrm{swap}_{04}\) is stated in ``dirac matrices'',
Secs. 0.1 and 2.1, and verified by the checks \texttt{71-N1} to
\texttt{71-N3}; the registry marks the complex set with the word NOT.

\textbf{N03, \texttt{Bessel\ product.nb}} (title: ``Bessel Product'').
Identities for products of Bessel functions of arguments proportional to
\(e^{\pm t/6}\), and coefficients of an expansion in such products. It
defines six symbols (\texttt{cApprox}, \texttt{capprox}, \texttt{rc},
\texttt{rC}, \texttt{eqΦ0007}, \texttt{sν}), none of which occurs in
another notebook or in the project. Not used.

\textbf{N04, \texttt{Dirac\_Einstein-Lovelock\_4+4\_Nash.nb}} (title:
``Scratch work of Author''). Gravity on the 4+4 spacetime with the
tensor package xAct. \emph{Conventions:}
\texttt{X\ =\ \{x0{[}{]},\ …,\ x7{[}{]}\}}; its text states that \(x0\)
is a hidden space coordinate (a small circle), \(x1,x2,x3\) the usual
space coordinates, \(x4\) the time and \(x5,x6,x7\) ``superluminal
deflating time'' coordinates. \emph{Symbols taken over:} the list
\texttt{X} and the meaning of the coordinates \(\to X=\{x0,\dots,x7\}\)
of the project. \emph{Collisions:} \texttt{g44} is the metric tensor of
a curved 4+4 spacetime (the first version of the project's programs used
this name for the flat Dirac matrices: error E7); \texttt{H} is a
fundamental inverse length of the field equations and \texttt{Λ} their
cosmological term (in the project \(\mathsf h\) is a Hamiltonian matrix,
\(H_0\) the Hubble constant and \(\Lambda\) a Lorentz matrix or a
cutoff); \texttt{β1}, \texttt{β2}, \texttt{β3} are functions in the
metric (in the project \(\beta\) is a matrix); the letters \texttt{a},
\texttt{b}, \texttt{c}, \texttt{d}, \texttt{f}, \texttt{i}, \texttt{j},
\texttt{k}, \texttt{l}, \texttt{m}, \texttt{n}, \texttt{p}, \texttt{q},
\texttt{r}, \texttt{s}, \texttt{y} are abstract tensor indices. The name
\texttt{T16\^{}A} occurs in a text cell only. \emph{Remedy:} the flat
metric of the project is \(\eta_{1644}\); \texttt{g44}, \texttt{H},
\texttt{Λ}, \texttt{β1} to \texttt{β3} are not used.

\textbf{N05, \texttt{Dirac\_Nash\_Identity.nb}.} Identities of real
\(4\times4\) matrices (the subject of the paper D06).
\emph{Conventions:} \texttt{X\ =\ \{x1,\ …,\ x6\}}, six coordinates.
\emph{Collisions:} \texttt{X} (six coordinates against eight);
\texttt{γ} (a list of four \(4\times4\) matrices); \texttt{σ} is the
\(4\times4\) matrix \texttt{γ64} here, not the \(8\times8\) matrix of
N01 (which N05 defines under the name \texttt{σ88} only); \texttt{η44}
is the \(4\times4\) metric diag\((1,1,1,-1)\); \texttt{S44αβ} denotes
commutators of the four matrices \texttt{γ} (nothing to do with
\texttt{S44} of N02). \texttt{η4488}, \texttt{σ88}, \texttt{s4by4},
\texttt{t4by4} are defined by the same statements as in N01.
\emph{Remedy:} nothing is taken from this notebook.

\textbf{N06, \texttt{Einstein-Lovelock-4+4-Nash.nb}; N07,
\texttt{Einstein-Lovelock-7+1-Nash.nb}; N14,
\texttt{field\_equations.nb}.} The Einstein--Lovelock field equations on
curved 4+4 and 7+1 spacetimes with xAct. \emph{Conventions:} the list
\texttt{X} of N04, with the same explanation of the coordinates.
\emph{Collisions:} \texttt{g44} (N06, N14) and \texttt{g71} (N07) are
curved metrics. In N06 and N07 a chapter title and a text cell that is
not evaluated still show the model metric in an older convention, with a
coordinate \texttt{x8} in the place of \texttt{x0} (the convention of
N08); the evaluated cells use \texttt{x0{[}{]}}. \texttt{H}, \texttt{Λ},
\texttt{β1} to \texttt{β3} and the index letters are as in N04.
\emph{Remedy:} nothing is taken from these notebooks; the second stage
of the author's programme, to which they belong, is not implemented in
this project.

\textbf{N08, \texttt{Einstein-Lovelock-Eternal.nb}.} A solution of the
same field equations. \emph{Conventions:} a chart with the coordinates
\texttt{x1{[}{]}} to \texttt{x8{[}{]}} and the index values 1 to 8; the
metric is called \texttt{g}; the coordinate \texttt{x8} occupies the
place of the hidden coordinate, and \texttt{β1}, \texttt{β2} are
functions of it. The cosmological term enters the vacuum equations with
the sign opposite to that of N04, N06, N07 and N14. \emph{Collision:}
the coordinate names are shifted with respect to N04 (\texttt{x8} here
corresponds to \texttt{x0} there; \texttt{x1} is the first coordinate
here and the second there). \emph{Remedy:} not used; the project's
coordinates are those of N04.

\textbf{N09, \texttt{Generalized\ \_Kronecker\_Delta\_4+4.nb}.} The
generalized Kronecker delta and the Lovelock tensors with xAct.
\emph{Conventions:} chart \(x0,\dots,x7\); metric \texttt{g} of
signature 4+4. \emph{Collisions:} besides the Latin index letters of
N04, the Greek letters \texttt{α}, \texttt{β}, \texttt{γ}, \texttt{δ},
\texttt{ε}, \texttt{ζ}, \texttt{η}, \texttt{κ}, \texttt{λ}, \texttt{μ},
\texttt{ν}, \texttt{ξ}, \texttt{ρ}, \texttt{σ}, \texttt{τ}, \texttt{χ}
are declared as abstract tensor indices; in N01 four of them
(\texttt{γ}, \texttt{ζ}, \texttt{σ}, \texttt{τ}) are matrices.
\texttt{g} is a metric here and in N08. \emph{Remedy:} not used.

\textbf{N10, N11, N12} (the three notebooks on the wave function of the
universe; their file names are in the table above; title of each:
``Scratch work of Author''). Work on a wave function of the universe on
the 4+4 spacetime; each repeats the construction of N01 with the same
statements, and N10 also contains the material of N05.
\emph{Conventions:} \texttt{X\ =\ \{x0,\ …,\ x7\}} as plain symbols,
with the meaning of N04. \emph{Collisions:} \texttt{Ψ16} is a list of
sixteen functions of \texttt{x0} and \texttt{x4}, called the wave
function of the universe; \texttt{ω} is the array of the coefficients of
a spin connection (in the project \(\omega_{1644}\) is a product of
Dirac matrices); \texttt{S} is a transformation of spinors, \texttt{Γ}
the Christoffel symbols; \texttt{SAB} and \texttt{sAB{[}A{]}{[}B{]}} are
the commutators \(\tfrac14[\)\texttt{T16\^{}A{[}A{]}}\(,\)
\texttt{T16\^{}A{[}B{]}}\(]\) (in N02 \texttt{sAB} is a sign);
\texttt{T16\^{}α} are Dirac matrices with a spacetime index; in the text
of N10, \(\pm M\) are masses of universes and \texttt{H} is ``some
function of M''. \emph{Remedy:} not used; the matrices are taken from
N01. The commutators \texttt{SAB} are the same matrices as
\(\Sigma_{1644}^{AB}\) of the project, which defines them independently.

\textbf{N13, \texttt{det-g.nb}.} A short calculation of a determinant
for a diagonal metric built from \texttt{β1}, \texttt{β2}, \texttt{β3}.
Not used.

\hypertarget{the-papers}{%
\subsection{5.2 The papers}\label{the-papers}}

\begin{longtable}[]{@{}
  >{\raggedright\arraybackslash}p{(\columnwidth - 4\tabcolsep) * \real{0.0385}}
  >{\raggedright\arraybackslash}p{(\columnwidth - 4\tabcolsep) * \real{0.3205}}
  >{\raggedright\arraybackslash}p{(\columnwidth - 4\tabcolsep) * \real{0.6410}}@{}}
\toprule\noalign{}
\begin{minipage}[b]{\linewidth}\raggedright
Id
\end{minipage} & \begin{minipage}[b]{\linewidth}\raggedright
File
\end{minipage} & \begin{minipage}[b]{\linewidth}\raggedright
Paper
\end{minipage} \\
\midrule\noalign{}
\endhead
\bottomrule\noalign{}
\endlastfoot
D01 & \texttt{6dim\ dirac.pdf} & E. A. Lord, ``The Dirac spinor in six
dimensions'', Proc. Camb. Phil. Soc. 64, 765 (1968) \\
D02 & \texttt{A\ Remarkable\ Representation.pdf} & P. A. M. Dirac, ``A
remarkable representation of the 3 + 2 de Sitter group'', J. Math. Phys.
4, 901 (1963) \\
D03 & \texttt{A\ new\ spin-half.pdf} & P. L. Nash, ``A new spin-1/2 wave
equation'', J. Math. Phys. 26, 1363 (1985) \\
D04 & \texttt{Bryant-octon-17-2-185-txt.pdf} & R. L. Bryant,
``Submanifolds and special structures on the octonians'', J.
Differential Geom. 17, 185 (1982) \\
D05 & \texttt{Calibrated\ geometries.pdf} & R. Harvey and H. B. Lawson,
Jr., ``Calibrated geometries'', Acta Math. 148, 47 (1982) \\
D06 & \texttt{Identities\ satisfied.pdf} & P. L. Nash, ``Identities
satisfied by the generators of the Dirac algebra'', J. Math. Phys. 25,
204 (1984) \\
D07 & \texttt{On\ the\ structure\ of\ the\ split\ octonion\ algebra.pdf}
& P. L. Nash, ``On the structure of the split octonion algebra'', Nuovo
Cimento B 105, 31 (1990) \\
D08 & \texttt{Second\ gravity.pdf} & P. L. Nash, ``Second gravity'', J.
Math. Phys. 51, 042501 (2010) \\
D09 & \texttt{Vector\ cross\ products.pdf} & R. B. Brown and A. Gray,
``Vector cross products'', Comment. Math. Helv. 42, 222 (1967) \\
\end{longtable}

\textbf{D07 and D08} define real \(8\times8\) matrices \(\tau^A\),
\(\bar\tau^A\) (there called reduced Brauer--Weyl generators) of the
kind from which the 4+4 Dirac matrices of N01 are built, and D08 defines
the \(16\times16\) matrices \(t^A\) with the blocks
\((0,\bar\tau^A;\,\tau^A,0)\). \emph{Conventions that differ from the
notebooks:} the index runs over \(A=1,\dots,8\), with
\(\tau^8=\bar\tau^8=1\); in N01 and in the project it runs over
\(A=0,\dots,7\), with \(\tau^0=\bar\tau^0=1\). The metric is called
\(G_{AB}\), with the diagonal \((1,1,1,-1,-1,-1,-1,1)\) for
\(A=1,\dots,8\). A tilde denotes the transpose (in the project a
superscript \(T\) does, and \(\tilde\psi\) is an adjoint). D07 adds
that, relative to an earlier paper of the same author, ``tau bar'' and
``tau'' are exchanged and the sign of \(\tau^7\) is changed; and it uses
the letter \(H\) for a product of the seven matrices with
\(A=1,\dots,7\), which it requires to be \(+1\) (in the scanned copy the
formula is not legible enough to say which of the factors carry a bar).
\emph{Remedy:} the project takes its matrices from N01, not from the
papers, and uses \(\tau_{1644}^A\), \(\bar\tau_{1644}^A\),
\(\gamma_{1644}^A\), \(\eta_{1644}\) with \(A=0,\dots,7\); a formula of
D07 or D08 with the index value 8 corresponds to the index value 0 of
the project. For the matrices of the project, \(\tau_{1644}^0=1\) and
\(\tau_{1644}^1\cdots\tau_{1644}^7=1\) (check \texttt{44-B1}). They were
not compared entry by entry with the matrices printed in the papers.

\textbf{D06} treats fifteen real \(4\times4\) matrices \(\gamma^{AB}\),
labelled by a pair of indices with the values 1 to 6 and a metric of
signature 3+3, and calls the six antisymmetric ones \(s^h\) (self-dual)
and \(t^h\) (anti-self-dual), \(h=1,2,3\). The notebooks N01 and N05
follow it: their \texttt{s4by4{[}h{]}}, \texttt{t4by4{[}h{]}} carry
these letters, and the list \texttt{γ} of N01 consists of sixteen
matrices of this kind (the six, their nine products
\texttt{s4by4{[}J{]}.t4by4{[}K{]}} and the unit matrix). \emph{Remedy:}
the project keeps the letters \(s_h\), \(t_h\) for the \(4\times4\)
blocks and never writes \(\gamma\) without two things: an upper index
with the values 0 to 7, and the tag of the spacetime where the matrix is
defined.

\textbf{D01} (Lord) gives the general procedure by which D07 and D08
build spinor representations from reduced generators; it works with the
Dirac matrices of four dimensions and with spinors of five and six
dimensions. \textbf{D03} derives a wave equation for a massive particle
of spin \(\tfrac12\) from a Casimir operator, and \textbf{D02}
constructs a representation of the 3+2 de Sitter group in terms of two
oscillators. In these three papers \(\gamma\), \(\sigma\), \(\psi\) and
the indices have four-, five- or six-dimensional meanings.
\emph{Remedy:} none of their symbols is used.

\textbf{D04, D05 and D09} (Bryant; Harvey and Lawson; Brown and Gray)
are the mathematical literature on the octonions, on which the 7+1
construction rests. The project quotes them for the existence and the
properties of the octonion algebra and takes no formula in their
notation; its octonionic matrices are \(\tau_{1671}^A\), defined in its
own documents and verified by machine (checks \texttt{71-B1},
\texttt{71-D6}). The copies of these papers are scans, and the text
extracted from them does not render their Greek letters reliably; their
notation was therefore not tabulated. One collision is visible even so:
D09 writes \(X\) for a vector cross product, the letter that the
notebooks use for the list of coordinates.

\hypertarget{the-project-ncss}{%
\subsection{5.3 The project NCSS}\label{the-project-ncss}}

The task specification asks that the author's project NCSS (``Neutrino
Condensate Singlet Superfluid'', github.com/once-ere/NCSS, folder
\texttt{Revision\_1}) be read for context. Three things were taken from
it: the name Vorton-Scalar and its meaning; two data files with spectra
of the standard cosmological model, which serve as the reference in the
validation of the Boltzmann solver (\texttt{data/reference/}, with
checksums); and the scripts that build the documents and the
manuscripts, which were adapted (\texttt{build\_docs.sh},
\texttt{build\_manuscripts.py}, \texttt{tex2md.py}; each says so in its
header). NCSS has a symbol table of its own (its Supplement S4). The
symbols that matter here:

\begin{longtable}[]{@{}
  >{\raggedright\arraybackslash}p{(\columnwidth - 4\tabcolsep) * \real{0.1368}}
  >{\raggedright\arraybackslash}p{(\columnwidth - 4\tabcolsep) * \real{0.4105}}
  >{\raggedright\arraybackslash}p{(\columnwidth - 4\tabcolsep) * \real{0.4526}}@{}}
\toprule\noalign{}
\begin{minipage}[b]{\linewidth}\raggedright
NCSS
\end{minipage} & \begin{minipage}[b]{\linewidth}\raggedright
Meaning in NCSS
\end{minipage} & \begin{minipage}[b]{\linewidth}\raggedright
In this project
\end{minipage} \\
\midrule\noalign{}
\endhead
\bottomrule\noalign{}
\endlastfoot
\(\chi_V\), also \(\chi\) & the Vorton-Scalar: a \emph{real} scalar
field & \(\Phi_{VS}\), a \emph{complex} scalar field; \(\chi\) is the
chirality matrix \\
\(m_\chi\) & mass of the Vorton-Scalar & \(m_s\), mass of the radial
mode; the phase mode \(\vartheta_{VS}\) is massless \\
\(g\), \(g_{ij}\); \(\alpha=g^2/4\pi\) & coupling to the neutrinos (a
matrix in the space of the three species) & \(g_{VS}=m/v_{VS}\);
\(\alpha_{VS}=g_{VS}^2/4\pi\) \\
\(D\), \(D_c\) & \(D=\alpha m_\nu/2m_\chi\); two neutrinos bind for
\(D>D_c=0.83990\) & \(D_{VS}=\alpha_{VS}m/m_s\), \textbf{twice} the
combination of NCSS; two fermions bind for \(D_{VS}>D_c=1.67981\) \\
\(C\) & \(g^2m_\nu^2/m_\chi^2\), the ``nugget parameter''; also the
charge-conjugation matrix of four dimensions & \(C=g_{VS}^2m^2/m_s^2\),
the same combination for the fields of this project; the matrices are
\(C_{1644}^{\mp}\), \(C_{1671}^{\pm}\) \\
\(\phi\) & \(\langle\chi_V\rangle\) & \(\phi\) with
\(\langle\Phi_{VS}\rangle=\phi/\sqrt2\) \\
\(\Delta_{ij}\) & gap matrix in the space of the three neutrino species
& \(\Delta^a\), \(a\in\{0,5,6,7\}\), a vector along the four extra
axes \\
\(\nu_i\) & the neutrino fields & \(\psi_{1644}\), \(\psi_{1671}\);
\(\nu\) is the number 8 of internal states \\
\(S\) & a stability parameter of neutrino matter & Sec. 4 \\
\(\lambda\) & \((D/\pi x_F)\ln(1+4x_F^2)\), with \(x_F=k_F/m_\chi\) &
\(\lambda=N(0)\,W(k_F,k_F)\); \(x\) is \(m^*/m\) \\
metric & four dimensions, signature \((+,-,-,-)\) & eight dimensions:
\(\eta_{1644}\) (space-like \(=+1\)), \(\eta_{1671}\) (time-like
\(=+1\)) \\
\end{longtable}

Two warnings follow from the table. A value of \(D\) must not be carried
from one project to the other without the factor 2. And the thresholds
of the coupling \(C\) quoted in the two projects are not comparable: in
NCSS the scalar is real with a mass term and each species has two
internal states, here the scalar is complex with a symmetry-breaking
potential and the field has eight internal states (Supplement S4, Sec.
5). \emph{Remedy:} every quantity of the Vorton-Scalar of this project
carries the suffix VS in the registry, so that it cannot be mistaken for
the quantity of NCSS with the same letter.

\begin{center}\rule{0.5\linewidth}{0.5pt}\end{center}

\hypertarget{collisions-between-the-original-files-symbol-by-symbol}{%
\section{6. Collisions between the original files, symbol by
symbol}\label{collisions-between-the-original-files-symbol-by-symbol}}

\hypertarget{the-same-symbol-for-different-objects}{%
\subsection{6.1 The same symbol for different
objects}\label{the-same-symbol-for-different-objects}}

The table lists the symbols of the original files that (i) have
different meanings in different files and (ii) concern an object that
the project uses. The complete machine-made list of the symbols that
several notebooks define is in Appendix A.

\begin{longtable}[]{@{}
  >{\raggedright\arraybackslash}p{(\columnwidth - 4\tabcolsep) * \real{0.1205}}
  >{\raggedright\arraybackslash}p{(\columnwidth - 4\tabcolsep) * \real{0.5783}}
  >{\raggedright\arraybackslash}p{(\columnwidth - 4\tabcolsep) * \real{0.3012}}@{}}
\toprule\noalign{}
\begin{minipage}[b]{\linewidth}\raggedright
Symbol
\end{minipage} & \begin{minipage}[b]{\linewidth}\raggedright
Meaning, by file
\end{minipage} & \begin{minipage}[b]{\linewidth}\raggedright
In the project
\end{minipage} \\
\midrule\noalign{}
\endhead
\bottomrule\noalign{}
\endlastfoot
\texttt{X}, \texttt{x0} \ldots{} \texttt{x7} & N01, N04, N06, N07,
N10--N12: coordinates \(x0,\dots,x7\), with the statement that \(x0\) is
a hidden space coordinate and \(x4\) the time. N09, N14: the same chart
without explanation. N02: the same names, \(x0\) time. N05: \texttt{X}
\(=\{x1,\dots,x6\}\). N08: coordinates \(x1,\dots,x8\), \(x8\) hidden &
\(X=\{x0,\dots,x7\}\) with the meaning of N01 and N04, never renamed;
N02 by \(\mathrm{swap}_{04}\) \\
\texttt{x8} & N08: the hidden coordinate. N06, N07: in a chapter title
and a text cell, where the evaluated cells have \texttt{x0} & not
used \\
\texttt{g44}, \texttt{g71}, \texttt{g} & N04, N06, N14: \texttt{g44},
curved metric of signature 4+4. N07 and N02 (Sec. 12.2): \texttt{g71},
curved metric of signature 7+1. N08, N09: \texttt{g}, curved metric &
the flat metrics are \(\eta_{1644}\), \(\eta_{1671}\); \texttt{g44},
\texttt{g71} never denote project objects \\
\texttt{η17}, \texttt{η71} & N02: two metrics of the 7+1 signature,
\texttt{η17} \(=-\)\texttt{η71} & \(\eta_{1671}\) is \texttt{η17} with
the index values 0 and 4 exchanged \\
\texttt{η44} & N05, N10: the \(4\times4\) metric diag\((1,1,1,-1)\) (in
N05 through a symbol \texttt{g\_αβ}) & not used; the error E7 concerned
names of this kind \\
\texttt{γ} & N01, N10--N12: sixteen \(4\times4\) matrices. N05: four
\(4\times4\) matrices. N09: tensor index. D06: \(\gamma^{AB}\), fifteen
\(4\times4\) matrices & \(\gamma_{1644}^A\), \(\gamma_{1671}^A\): real
\(16\times16\) Dirac matrices \\
\texttt{Γ17}, \texttt{Γ71}, \texttt{Γ71R32}; \texttt{Γ} & N02: real
\(16\times16\), complex \(16\times16\) and real \(32\times32\)
generators of 7+1 Clifford algebras. N10--N12: \texttt{Γ}, Christoffel
symbols & \(\gamma_{1671}^A\) from \texttt{Γ17up}; the other two sets
are not used \\
\texttt{τ}, \texttt{ζ}, \texttt{τO} & N01, N02, N10--N12:
\texttt{τ{[}A{]}} (and \texttt{ζ{[}A{]}}), \(A=0,\dots,7\), the
\(8\times8\) blocks of 4+4. N02: \texttt{τO{[}A{]}}, the octonionic
matrices. N09: \texttt{τ}, \texttt{ζ} tensor indices. D07, D08:
\(\tau^A\), \(A=1,\dots,8\) & \(\tau_{1644}^A\) and \(\tau_{1671}^A\);
the plain letter \(\tau\) is the optical depth \\
\texttt{σ}, \texttt{σ8} & N01, N02, N10--N12: \texttt{σ}, the
\(8\times8\) matrix \((0,1_4;1_4,0)\). N05: \texttt{σ} \(=\) \texttt{σ8}
\(=\) \texttt{γ64}, a \(4\times4\) matrix. N09: tensor index. D08:
\(\sigma\) lowers spinor indices. N02 (Secs. 7.11, 9.1): \(\sigma^A\),
\(\bar\sigma^A\), a family of \(8\times8\) matrices & \(\sigma_{88}\);
other uses of \(\sigma\) in Sec. 4 \\
\texttt{S44}, \texttt{S44αβ} & N02: \texttt{S44}, the product of the
four time-like 4+4 Dirac matrices. N05, N10: \texttt{S44αβ}, commutators
of four \(4\times4\) matrices & \(C_{1644}^{+}\) \\
\texttt{sAB}, \texttt{SAB} & N10--N12: the commutators
\(\tfrac14[\)\texttt{T16\^{}A{[}A{]}}\(,\)\texttt{T16\^{}A{[}B{]}}\(]\).
N02: \texttt{sAB}, a sign & \(\Sigma_{1644}^{AB}\) \\
\texttt{ω}, \texttt{ω17} & N10--N12: \texttt{ω}, coefficients of a spin
connection. N02: \texttt{ω17} \(=\) \texttt{Γ17{[}8{]}}; in its text
\(\omega\) is also a frequency & \(\omega_{1644}\), \(\omega_{1671}\):
products of the eight Dirac matrices \\
\texttt{χ17}, \texttt{χ71}; \texttt{χ} & N02: chirality matrices of the
real and of the complex set; in the text of its Sec. 13.2, \(\chi\) is a
rescaled field. N09: \texttt{χ} tensor index. NCSS: \(\chi\) is the
Vorton-Scalar & \(\chi_{1671}\) from \texttt{χ17} \\
\texttt{P\_L}, \texttt{P\_R}, \texttt{PL…}, \texttt{PR…} & N01,
N10--N12: \texttt{P\_L}, \texttt{P\_R} (4+4). N02: \texttt{PL44},
\texttt{PR44} (4+4); \texttt{PL17}, \texttt{PR17} (real 7+1 set);
\texttt{PL71}, \texttt{PR71} (complex set) & \(P_{1644}^{L,R}\),
\(P_{1671}^{L,R}\) \\
\texttt{Ψ16}, \texttt{ψ}, \texttt{Ψ} & N10--N12: \texttt{Ψ16}, sixteen
functions of \texttt{x0}, \texttt{x4} (wave function of the universe).
N02: \(\psi\), the field \texttt{dirac16complex}; \texttt{Ψ}, a pair of
octonions & \(\psi_{1644}\), \(\psi_{1671}\) (dirac1644, dirac1671) \\
\texttt{H} & N04, N06, N07, N08, N14 (and N02, Sec. 12.2): a fundamental
inverse length, declared as a constant. N02: Hamiltonian (Sec. 7.7),
flavour Hamiltonian (Sec. 19). N10: a function of a mass \(M\). D07: the
product \(\tau^1\cdots\tau^7\) & \(\mathsf h_{1644}\),
\(\mathsf h_{1671}\) for the Hamiltonian matrices; \(H_0\) for the
Hubble constant \\
\texttt{Λ} & N04, N06, N07, N08, N14: cosmological term (with opposite
signs in N08 and in the others). N02: Lorentz matrix (Sec. 5.1), cutoff
(Sec. 9.6) & Lorentz matrix or cutoff (Sec. 4) \\
\texttt{β1}, \texttt{β2}, \texttt{β3} & N04, N06, N07, N13, N14:
functions of \texttt{x0{[}{]}} and \texttt{x4{[}{]}} in a metric. N08:
functions of \texttt{x8{[}{]}} and \texttt{x4{[}{]}}. N10--N12: the same
with plain \texttt{x0}, \texttt{x4} & not used; \(\beta_{1644}\),
\(\beta_{1671}\) are matrices \\
index letters & N02, N04, N06--N09, N14: \texttt{a}, \texttt{b},
\texttt{c}, \texttt{d}, \texttt{f}, \texttt{i}, \texttt{j}, \texttt{k},
\texttt{l}, \texttt{m}, \texttt{n}, \texttt{p}, \texttt{q}, \texttt{r},
\texttt{s}, \texttt{y} are abstract tensor indices; N09 adds sixteen
Greek letters. D07, D08: \(A=1,\dots,8\). D06: \(A,B=1,\dots,6\);
\(h=1,2,3\) & \(A,B=0,\dots,7\); \(a,b\in\{0,5,6,7\}\); \(i,j\)
space-like; \(t=5,6,7\); \(h=1,2,3\) (Sec. 3.1) \\
\end{longtable}

\hypertarget{different-symbols-for-the-same-object}{%
\subsection{6.2 Different symbols for the same
object}\label{different-symbols-for-the-same-object}}

The converse also occurs: the same object carries different names in
different files. The dictionary of \texttt{7+1.nb} (its Sec. 1.3) lists
the renamings that N02 made; the table adds the papers and the names of
the project.

\begin{longtable}[]{@{}
  >{\raggedright\arraybackslash}p{(\columnwidth - 8\tabcolsep) * \real{0.2836}}
  >{\raggedright\arraybackslash}p{(\columnwidth - 8\tabcolsep) * \real{0.1791}}
  >{\raggedright\arraybackslash}p{(\columnwidth - 8\tabcolsep) * \real{0.1791}}
  >{\raggedright\arraybackslash}p{(\columnwidth - 8\tabcolsep) * \real{0.1791}}
  >{\raggedright\arraybackslash}p{(\columnwidth - 8\tabcolsep) * \real{0.1791}}@{}}
\toprule\noalign{}
\begin{minipage}[b]{\linewidth}\raggedright
Object
\end{minipage} & \begin{minipage}[b]{\linewidth}\raggedright
N01 (and N10--N12)
\end{minipage} & \begin{minipage}[b]{\linewidth}\raggedright
N02
\end{minipage} & \begin{minipage}[b]{\linewidth}\raggedright
Papers
\end{minipage} & \begin{minipage}[b]{\linewidth}\raggedright
Project
\end{minipage} \\
\midrule\noalign{}
\endhead
\bottomrule\noalign{}
\endlastfoot
\(8\times8\) blocks & \texttt{τ{[}A{]}}, \texttt{τ} with an overbar &
\texttt{τ{[}A{]}}, \texttt{τbar{[}A{]}} & D07, D08: \(\tau^A\),
\(\bar\tau^A\) (\(A\) shifted) & \(\tau_{1644}^A\),
\(\bar\tau_{1644}^A\) \\
4+4 Dirac matrices & \texttt{T16\^{}A{[}A{]}} & \texttt{T16{[}A{]}} &
D08: \(t^A\) & \(\gamma_{1644}^A\) \\
their product & \texttt{T16\^{}A{[}8{]}} & \texttt{T16{[}8{]}} & &
\(\omega_{1644}\) \\
chiral projectors of 4+4 & \texttt{P\_L}, \texttt{P\_R} & \texttt{PL44},
\texttt{PR44} & & \(P_{1644}^{L,R}\) \\
metric of 4+4 & \texttt{η4488}, \texttt{η8}, \texttt{g\_AB} &
\texttt{η4488}, \texttt{gAB44} & D08: \(G_{AB}\) (\(A\) shifted) &
\(\eta_{1644}\) \\
product of the four space-like matrices & \texttt{σ16} & \texttt{σ16} &
& \(C_{1644}^{-}\) \\
product of the four time-like matrices & & \texttt{S44} & &
\(C_{1644}^{+}\) \\
commutators & \texttt{SAB}, \texttt{sAB{[}A{]}{[}B{]}} (N10--N12) & & &
\(\Sigma_{1644}^{AB}\) \\
\(4\times4\) blocks & \texttt{s4by4{[}h{]}}, \texttt{t4by4{[}h{]}} &
\texttt{s4by4{[}h{]}}, \texttt{t4by4{[}h{]}} & D06: \(s^h\), \(t^h\) &
\(s_h\), \(t_h\) \\
the \(8\times8\) matrix \((0,1_4;1_4,0)\) & \texttt{σ}, \texttt{σ88} &
\texttt{σ} & D08: \(\sigma\) & \(\sigma_{88}\) \\
\end{longtable}

\hypertarget{plain-letters-of-the-project-that-coincide-with-a-symbol-of-a-notebook}{%
\subsection{6.3 Plain letters of the project that coincide with a symbol
of a
notebook}\label{plain-letters-of-the-project-that-coincide-with-a-symbol-of-a-notebook}}

A tagged symbol cannot coincide with a symbol of a notebook (check
\texttt{S6}). A plain letter can. The following table lists every plain
symbol of the registry of Sec. 3 that is, as a string of letters, also a
symbol of a notebook (33 symbols), with both meanings. It is generated
by the script that performs the checks, and check \texttt{S9} fails if a
coincidence is missing from it. The section numbers are those of
\texttt{7+1.nb} (N02).

\begin{longtable}[]{@{}
  >{\raggedright\arraybackslash}p{(\columnwidth - 4\tabcolsep) * \real{0.0795}}
  >{\raggedright\arraybackslash}p{(\columnwidth - 4\tabcolsep) * \real{0.5000}}
  >{\raggedright\arraybackslash}p{(\columnwidth - 4\tabcolsep) * \real{0.4205}}@{}}
\toprule\noalign{}
\begin{minipage}[b]{\linewidth}\raggedright
Symbol
\end{minipage} & \begin{minipage}[b]{\linewidth}\raggedright
In the notebooks
\end{minipage} & \begin{minipage}[b]{\linewidth}\raggedright
In this project
\end{minipage} \\
\midrule\noalign{}
\endhead
\bottomrule\noalign{}
\endlastfoot
\texttt{a} & N02, N04, N06--N09, N14: abstract tensor index of the
manifold \texttt{M8} (package xAct) & index of the four extra axes,
values 0, 5, 6, 7 \\
\texttt{B} & N01, N10: \texttt{B\ =\ 1}, an index value in a check; N02:
index with the values 0 to 7 & index with the values 0 to 7; bag
constant (S4, M IX); Killing form \(B(X,Y)\) (DM, M II) \\
\texttt{b} & as \texttt{a} & as \(a\) \\
\texttt{C} & N02: index with the values 0 to 7 & index with the values 0
to 7; the coupling \(C=g_{VS}^2m^2/m_s^2\) \\
\texttt{h} & N01, N10: \texttt{h\ =\ hUSE}, the number 8 that selects a
unit spinor; N02: Hamiltonian (Sec. 7.3), flavour Hamiltonian (Sec. 19),
induced metric (Secs. 20.3, 25.6) & index of the matrices \(s_h\),
\(t_h\); Hubble constant in units of 100 km s\(^{-1}\) Mpc\(^{-1}\) (the
Hamiltonian matrix is \(\mathsf h\)) \\
\texttt{i} & N02, N04, N06--N09, N14: abstract tensor index (xAct) &
index of space-like coordinates \\
\texttt{j} & N02, N04, N06--N09, N14: abstract tensor index (xAct); text
of N10--N12: the number 134 of a basis matrix & index of space-like
coordinates \\
\texttt{M} & N02: an intertwining matrix (Sec. 4.4), the \(8\times8\)
matrix \(m+p_i\,\)\texttt{τO{[}i{]}} (Sec. 7.4), a Kohn--Sham mass (Sec.
9.3) & mass of an exchanged boson (the matrix of N02, Sec. 7.4, is
\(M_{1671}\); the gap matrix is \(\mathsf M\)) \\
\texttt{m} & N02, N04, N06--N09, N14: abstract tensor index (xAct); in
the text of N02 also the mass of the field & mass of the fermion field,
\(m_{1644}\) or \(m_{1671}\) with the tag dropped \\
\texttt{P} & N02: the \(32\times32\) matrix
\(\sum_Ap^A\,\)\texttt{Γ71R32{[}A{]}} (Sec. 5.5), a pseudoscalar
potential (Sec. 8.4) & generator of \(U(1)_E\); pressure; thermal boson
loop \(P(M)\) \\
\texttt{Q} & N02 (Secs. 1.2, 1.4): the quadratic form of a Clifford
algebra & thermal part \(Q(M)\) of a boson loop; electric charge (M
X) \\
\texttt{S} & N02: a unitary matrix (Sec. 4.3), a scalar potential (Sec.
8.4), the action (Secs. 12, 17); N10--N12: the transformation
\(\exp(\tfrac12\omega\,\)\texttt{SAB}\()\) of spinors & element of the
spin group (the meaning of N10--N12); classical action \(S[\Phi]\);
entropy; in S1 and S4 also \(\sum_a\lvert\Delta^a\rvert^2\) \\
\texttt{U} & N02 (Sec. 20.2): a stabilizing potential & the matrix
\(U(\vec p)\) of positive-energy spinors \\
\texttt{V} & N02 (Secs. 8.4, 9.6): a vector potential of the Kohn--Sham
theory & effective potential \(V[\Delta]\) \\
\texttt{W} & N02: the unitary matrix (\texttt{EL}, \texttt{ER}) (Sec.
7.11); a function \(a^3b^4U\) (Sec. 20.2) & kernel \(W(k,k')\) of the
gap equation; generating functional \(W[J,K]\) (the grading is
\(W_{1644}\)) \\
\texttt{x} & N02 (Sec. 1.2): an octonion \(x=x0\,e_0+\dots+x7\,e_7\) &
the ratio \(m^*/m\) \\
\texttt{β} & N09: abstract tensor index; N02: rapidity of a boost (Sec.
7.9) & the matrix \(\beta_{1644}\) or \(\beta_{1671}\) with the tag
dropped \\
\texttt{Γ} & N10--N12: the Christoffel symbols returned by a function
\texttt{rt} & effective action \(\Gamma[\phi,G]\); in M VI and S1 a
vertex matrix \\
\texttt{ε} & N09: abstract tensor index; N02: the matrix \((0,1;-1,0)\)
(Sec. 4), an energy density (Sec. 8.7) & energy density
\(\varepsilon(k_F,x)\) \\
\texttt{η} & N09: abstract tensor index; N02: short for \texttt{η17} &
eigenvalue \(\eta(T)\) of the linearized gap equation (the metrics carry
a tag or indices) \\
\texttt{κ} & N09: abstract tensor index; N02: a chiral density (Sec.
9.1), the constant \(8\pi G_8\) (Sec. 11) & pair correlation
\(\kappa(k)\) \\
\texttt{Λ} & N04, N06--N08, N14: cosmological term of the field
equations (xAct constant); N02: Lorentz matrix (Sec. 5.1), momentum
cutoff (Sec. 9.6) & Lorentz matrix \(\Lambda^A{}_B\); momentum cutoff \\
\texttt{λ} & N09: abstract tensor index; N02: coordinate index &
dimensionless pairing coupling \\
\texttt{μ} & N09: abstract tensor index; N02: coordinate index; chemical
potential (Sec. 8.7) & chemical potential \\
\texttt{ν} & N09: abstract tensor index; N02: coordinate index & the
number of internal states, \(\nu=8\) \\
\texttt{σ} & N01, N02, N10--N12: the \(8\times8\) matrix that is
\(\sigma_{88}\) here; N05: the \(4\times4\) matrix \texttt{γ64}; N09:
abstract tensor index & self-energy \(\sigma(p)\) (the matrix is always
written \(\sigma_{88}\)) \\
\texttt{τ} & N01, N02, N10--N12: the matrices \texttt{τ{[}A{]}}; N09:
abstract tensor index & optical depth of reionization (the matrices are
always written \(\tau_{1644}^A\), \(\tau_{1671}^A\)) \\
\texttt{φ} & N02 (Sec. 20.3): the canonical radion & expectation value
of the Vorton-Scalar \\
\texttt{χ} & N09: abstract tensor index; N02 (Sec. 13.2): a rescaled
field & chirality matrix of the spacetime at hand \\
\texttt{ψ} & N02: the field \texttt{dirac16complex} in formulas & the
fermion field (\(\psi_{1644}\) or \(\psi_{1671}\) where it is
defined) \\
\texttt{E\_x} & N02 (Sec. 8.4): exchange energy of its Kohn--Sham theory
& exchange (Fock) energy of the two modes of the Vorton-Scalar: the same
kind of quantity in another model \\
\texttt{m\_s} & N02 (Sec. 17.4): mass of the strange quark & mass of the
radial mode of the Vorton-Scalar (name in files: \texttt{msVS}) \\
\texttt{λ\_P} & N02 (Sec. 24.3): coupling of a Higgs portal & the part
of the pairing coupling that comes from the phase mode \\
\end{longtable}

In none of these cases does the project take the object of the notebook
under the plain letter: where it takes the object (the matrices
\texttt{σ}, \texttt{τ{[}A{]}}, the field), it uses the tagged symbol.
Two coincidences deserve attention because the letters carry a subscript
and look specific: \(m_s\) is the mass of the strange quark in N02 and
the mass of the radial mode of the Vorton-Scalar here (name in files:
\texttt{msVS}); and \(\lambda_P\) is a coupling of a Higgs portal in N02
and a part of the pairing coupling here.

\begin{center}\rule{0.5\linewidth}{0.5pt}\end{center}

\hypertarget{collisions-that-the-project-produced-itself-and-their-removal}{%
\section{7. Collisions that the project produced itself, and their
removal}\label{collisions-that-the-project-produced-itself-and-their-removal}}

The project is not exempt from the problem it is asked to solve. Four
collisions or errors of this kind were made and removed; each is in the
error log (\texttt{provenance/P9\_error\_log}).

\begin{enumerate}
\def\labelenumi{\arabic{enumi}.}
\tightlist
\item
  \textbf{E7.} The first version of the programs called the project's
  matrices \texttt{g44}, \texttt{g71}, \texttt{eta44}, \texttt{eta71},
  \texttt{gamma44}, \texttt{gamma71}, \texttt{chi71}. The notebooks use
  \texttt{g44} and \texttt{g71} for curved metrics, and \texttt{η71},
  \texttt{Γ71}, \texttt{χ71} for the \emph{complex} 7+1 set. All names
  were replaced by the tagged names of Sec. 3; check \texttt{S8}
  verifies that the old names survive only where a notebook is quoted.
\item
  \textbf{E35.} The output of the program used the letters A and B both
  for the two factors of the internal symmetry of the zero-mode sector
  and for the two factors of the rotation group of the four space
  coordinates of 4+4. These are different groups. The second pair is now
  called \(SU(2)_X\), \(SU(2)_Y\).
\item
  \textbf{E41.} The text of check \texttt{44-D6} defined
  \texttt{alpha1644\^{}a} as \(\gamma_{1644}^4\gamma_{1644}^a\), with
  the sign opposite to the definition
  \(\alpha_{1644}^j=-\gamma_{1644}^4\gamma_{1644}^j\) of S1, of the
  manuscripts and of the program itself (\texttt{kernel.rs}). The
  statement that the check verifies is true for either sign; the text
  was corrected.
\item
  \textbf{E42.} The first draft of this provenance was written partly
  from working notes, and the first version of the program that scans
  the notebooks recorded only definitions at the beginning of a line.
  Checking every statement against the notebooks showed that both were
  wrong in several places; the scanner was rewritten and the statements
  corrected before anything was committed.
\end{enumerate}

Section 4 records seven letters that were changed in the documents for
the same reason.

\begin{center}\rule{0.5\linewidth}{0.5pt}\end{center}

\hypertarget{the-machine-checks}{%
\section{8. The machine checks}\label{the-machine-checks}}

\begin{longtable}[]{@{}
  >{\raggedright\arraybackslash}p{(\columnwidth - 4\tabcolsep) * \real{0.0833}}
  >{\raggedright\arraybackslash}p{(\columnwidth - 4\tabcolsep) * \real{0.8194}}
  >{\raggedright\arraybackslash}p{(\columnwidth - 4\tabcolsep) * \real{0.0972}}@{}}
\toprule\noalign{}
\begin{minipage}[b]{\linewidth}\raggedright
Check
\end{minipage} & \begin{minipage}[b]{\linewidth}\raggedright
Statement
\end{minipage} & \begin{minipage}[b]{\linewidth}\raggedright
Result
\end{minipage} \\
\midrule\noalign{}
\endhead
\bottomrule\noalign{}
\endlastfoot
\texttt{S1} & every name of the registry belongs to exactly one entry
(78 names in 113 entries) & passed \\
\texttt{S2} & no name of the registry is a symbol defined in any of the
14 notebooks (1921 notebook symbols compared; exempt: the 8 coordinate
names x0 \ldots{} x7, taken from N04 with the same meaning) & passed \\
\texttt{S3} & 65 of the 78 names of the registry carry the tag 1644 or
1671 or the suffix VS; of the other 13, 8 are the coordinate names
(deliberately those of the notebooks, S2) and 5 (F\_pseudoscalar,
F\_scalar, J5, sigma88, swap04) are symbols of no notebook & passed \\
\texttt{S4} & every name with the tag 1644 or 1671 that occurs in a data
file or a document is in the registry (48 different names in 86 files; 4
further tokens are names of columns, bibliography keys or program
variables) & passed \\
\texttt{S5} & every name with the suffix VS that occurs in a data file
or a document is in the registry (17 different names) & passed \\
\texttt{S6} & no symbol of any notebook contains the tag 1644 or 1671 or
ends with VS, so a tagged name cannot be mistaken for a symbol of a
notebook & passed \\
\texttt{S7} & the manuscripts contain none of the 34 original names of
the objects taken from the notebooks (T16, Γ17, η4488, χ17, τO, g44,
\ldots): they use the symbols of the registry only & passed \\
\texttt{S8} & the names g44, g71, eta44, eta71, gamma44, gamma71, chi71,
which the first version of the programs used for project matrices
although the notebooks use them for other objects (error E7), occur in
the programs and data files only where a notebook is quoted & passed \\
\texttt{S9} & 33 plain (untagged) symbols of the registry coincide, as
strings of letters, with a symbol of a notebook (B, C, E\_x, M, P, Q, S,
U, V, W, a, b, h, i, j, m, m\_s, x, Γ, Λ, β, ε, η, κ, λ, λ\_P, μ, ν, σ,
τ, φ, χ, ψ); every one of them is listed with both meanings (table of
Sec. 6.3 of the MASTER provenance), and no listed coincidence is missing
& passed \\
\end{longtable}

The complete output is the file \texttt{data/symbol\_checks.txt}:

\begin{verbatim}
PROJECT spacetimes: checks of the symbol registry (provenance/symbols.csv) against the notebooks and the documents
9 of 9 checks passed

[PASS] S1  every name of the registry belongs to exactly one entry (78 names in 113 entries)
[PASS] S2  no name of the registry is a symbol defined in any of the 14 notebooks (1921 notebook symbols compared; exempt: the 8 coordinate names x0 ... x7, taken from N04 with the same meaning)
[PASS] S3  65 of the 78 names of the registry carry the tag 1644 or 1671 or the suffix VS; of the other 13, 8 are the coordinate names (deliberately those of the notebooks, S2) and 5 (F_pseudoscalar, F_scalar, J5, sigma88, swap04) are symbols of no notebook
[PASS] S4  every name with the tag 1644 or 1671 that occurs in a data file or a document is in the registry (48 different names in 86 files; 4 further tokens are names of columns, bibliography keys or program variables)
[PASS] S5  every name with the suffix VS that occurs in a data file or a document is in the registry (17 different names)
[PASS] S6  no symbol of any notebook contains the tag 1644 or 1671 or ends with VS, so a tagged name cannot be mistaken for a symbol of a notebook
[PASS] S7  the manuscripts contain none of the 34 original names of the objects taken from the notebooks (T16, Γ17, η4488, χ17, τO, g44, ...): they use the symbols of the registry only
[PASS] S8  the names g44, g71, eta44, eta71, gamma44, gamma71, chi71, which the first version of the programs used for project matrices although the notebooks use them for other objects (error E7), occur in the programs and data files only where a notebook is quoted
[PASS] S9  33 plain (untagged) symbols of the registry coincide, as strings of letters, with a symbol of a notebook (B, C, E_x, M, P, Q, S, U, V, W, a, b, h, i, j, m, m_s, x, Γ, Λ, β, ε, η, κ, λ, λ_P, μ, ν, σ, τ, φ, χ, ψ); every one of them is listed with both meanings (table of Sec. 6.3 of the MASTER provenance), and no listed coincidence is missing
\end{verbatim}

What these checks do not prove: that every \emph{sentence} of the
documents uses the symbols correctly, and that the hand-written
statements of Secs. 4 to 6 are complete. The first is a matter of
reading; for the second, the two machine-made lists (Appendix A for the
notebooks, \texttt{data/symbol\_concordance.csv} for the documents of
the project) let a reader test any statement.

\begin{center}\rule{0.5\linewidth}{0.5pt}\end{center}

\hypertarget{appendix-a.-the-symbols-that-the-notebooks-define}{%
\section{Appendix A. The symbols that the notebooks
define}\label{appendix-a.-the-symbols-that-the-notebooks-define}}

\emph{Method.} Each notebook is converted to text, cell by cell, by
\texttt{tools/nb2txt.py}; typeset input appears in linear form, a
subscript as \texttt{P\_L}, a superscript as \texttt{T16\^{}A}, an
overbar as \texttt{Over{[}τ,\_{]}}. The script
\texttt{tools/scan\_notebook\_symbols.py} then records the symbols that
a notebook defines:

\begin{itemize}
\tightlist
\item
  in \textbf{input cells}, after comments and strings have been removed,
  the head of the left-hand side of \emph{every} assignment (\texttt{=},
  \texttt{:=}, \texttt{\^{}=}, \texttt{\^{}:=}), wherever it stands: at
  the start of a line, inside parentheses, inside \texttt{Do{[}...{]}}
  or \texttt{Table{[}...{]}}, or in a chain \texttt{a\ =\ b\ =\ c}.
  Local variables of \texttt{Module}, \texttt{With}, \texttt{Block} are
  not counted, nor are the limits of typeset sums, properties of
  built-in functions and settings whose names begin with \texttt{\$};
\item
  the \textbf{declarations of xAct}: the manifold and its abstract
  indices (\texttt{DefManifold}), the chart and its coordinates
  (\texttt{DefChart}), the metric and its covariant derivative
  (\texttt{DefMetric}), constants and scalar functions;
\item
  in \textbf{text cells and headings}, the symbol that stands
  immediately to the left of a single ``='' in a formula of the prose,
  because this is notation that a reader of the notebook meets. A symbol
  that a notebook introduces \emph{only} in this way is marked with an
  asterisk. This rule has limits, which the header of the script states:
  a product written without a blank is read as one name, and a bare word
  of four or more lower-case letters is taken as prose.
\end{itemize}

A notebook can define a symbol several times with different meanings;
the scan keeps the first definition (a definition in code before one in
prose), and Sec. 5 treats such cases by hand. Output and message cells
are ignored. The defining statements themselves are not written into the
repository, because the notebooks are unpublished; the option
\texttt{-\/-dump} of the script writes them to a file for checking.

\emph{Result.} 1921 symbols; 350 are defined in two or more notebooks.
The table groups these by the set of notebooks that define them (54
sets). ``Same defining text'' means that the first defining statement,
with blanks removed, is identical in all notebooks of the set; different
text does not prove that the objects differ, and identical text does not
prove that they agree (the right-hand sides may contain symbols that
differ). The column is a guide: 274 symbols have the same text, 76 have
not, and Secs. 5 and 6 treat every symbol that matters for the project
individually. The complete list, with the symbols that only one notebook
defines, is the file
\texttt{provenance/generated/notebook\_symbols.csv}.

\begin{longtable}[]{@{}
  >{\raggedright\arraybackslash}p{(\columnwidth - 4\tabcolsep) * \real{0.1429}}
  >{\raggedright\arraybackslash}p{(\columnwidth - 4\tabcolsep) * \real{0.4286}}
  >{\raggedright\arraybackslash}p{(\columnwidth - 4\tabcolsep) * \real{0.4286}}@{}}
\toprule\noalign{}
\begin{minipage}[b]{\linewidth}\raggedright
Notebooks
\end{minipage} & \begin{minipage}[b]{\linewidth}\raggedright
Symbols with the same defining text
\end{minipage} & \begin{minipage}[b]{\linewidth}\raggedright
Symbols with different defining text
\end{minipage} \\
\midrule\noalign{}
\endhead
\bottomrule\noalign{}
\endlastfoot
N02 N04 N06 N07 N08 N09 N10 N11 N12 N14 & & \texttt{j}*, \texttt{x4}*,
\texttt{x5}*, \texttt{x6}*, \texttt{x7}* \\
N02 N04 N05 N06 N07 N10 N11 N12 N14 & & \texttt{X} \\
N02 N04 N06 N07 N09 N10 N11 N12 N14 & & \texttt{x0}* \\
N04 N06 N07 N08 N09 N10 N11 N12 N14 & \texttt{dir}, \texttt{header},
\texttt{name}, \texttt{names} & \\
N04 N06 N07 N08 N10 N11 N12 N13 N14 & & \texttt{β1}, \texttt{β2},
\texttt{β3} \\
N02 N04 N06 N07 N08 N09 N13 N14 & \texttt{kδ} & \\
N02 N04 N06 N07 N10 N11 N12 N14 & & \texttt{MyArrayComponents}* \\
N04 N06 N07 N08 N09 N10 N11 N12 & \texttt{time1} & \\
N04 N06 N07 N08 N10 N11 N12 N14 & \texttt{A\_a\^{}b}*, \texttt{DIM8}
& \\
N01 N02 N05 N09 N10 N11 N12 & & \texttt{γ}*, \texttt{σ} \\
N02 N04 N06 N07 N08 N09 N14 & \texttt{CD}, \texttt{M8}, \texttt{a},
\texttt{b}, \texttt{c}, \texttt{chartM8}, \texttt{d}, \texttt{f},
\texttt{i}, \texttt{k}, \texttt{l}, \texttt{m}, \texttt{n}, \texttt{p},
\texttt{q}, \texttt{r}, \texttt{s}, \texttt{y} & \texttt{x1},
\texttt{x2}, \texttt{x3} \\
N02 N04 N06 N07 N08 N10 N14 & & \texttt{H}* \\
N04 N06 N07 N10 N11 N12 N14 & \texttt{sX0}, \texttt{ssX} & \\
N01 N02 N05 N10 N11 N12 & \texttt{ID4}, \texttt{ID8}, \texttt{Qa},
\texttt{Qb}, \texttt{SelfDualAntiSymmetric}, \texttt{Zero4},
\texttt{Zero8}, \texttt{evalues}, \texttt{evecs}, \texttt{s4by4},
\texttt{t4by4} & \texttt{AntiSelfDualAntiSymmetric}, \texttt{η4488} \\
N01 N02 N09 N10 N11 N12 & & \texttt{τ} \\
N02 N04 N06 N07 N08 N14 & & \texttt{Λ}* \\
N04 N06 N07 N08 N09 N14 & \texttt{dimM8} & \texttt{indices} \\
N01 N02 N09 N11 N12 & & \texttt{ζ} \\
N01 N02 N10 N11 N12 & \texttt{Zero16} & \texttt{ID16}, \texttt{Ω}*,
\texttt{σ16} \\
N01 N04 N10 N11 N12 & & \texttt{T16\^{}A}* \\
N01 N05 N10 N11 N12 & \texttt{ε3}, \texttt{ε4} & \\
N02 N04 N06 N07 N08 & \texttt{EinsteinCD}, \texttt{RicciCD},
\texttt{RicciScalarCD}, \texttt{RiemannCD} & \\
N02 N05 N10 N11 N12 & & \texttt{u}* \\
N02 N06 N07 N08 N09 & & \texttt{Lovelock1}*, \texttt{Lovelock2}*,
\texttt{Lovelock3}* \\
N04 N06 N07 N08 N14 & \texttt{org} &
\texttt{EinsteinLovelockVacuumFieldEquations}, \texttt{allLovelock1},
\texttt{allLovelock2}, \texttt{allLovelock3} \\
N04 N06 N10 N11 N12 & & \texttt{MatrixMetric44} \\
N01 N10 N11 N12 & \texttt{BASE16}, \texttt{G16}, \texttt{Id16},
\texttt{MId}, \texttt{MId16}, \texttt{Over{[}τ,\_{]}}, \texttt{P\_L},
\texttt{P\_R}, \texttt{Simp}, \texttt{Simp16}, \texttt{ZERO16},
\texttt{anti}, \texttt{bas64}, \texttt{base16},
\texttt{determineIF8AntiSymmetric}, \texttt{dups}, \texttt{dupssymm16},
\texttt{find8AntiSymmetric}, \texttt{mid}, \texttt{simp}, \texttt{st},
\texttt{symm}, \texttt{symm16}, \texttt{t160}, \texttt{t16A},
\texttt{t16AB}, \texttt{t16ABC}, \texttt{t16ABCD}, \texttt{t16ABCDE},
\texttt{t16ABCDEF}, \texttt{t16ABCDEFG}, \texttt{t16ABCDEFGH},
\texttt{tA}, \texttt{tAB}, \texttt{tABC}, \texttt{use16}, \texttt{η64},
\texttt{σΩ} & \\
N02 N10 N11 N12 & & \texttt{S}*, \texttt{sAB}, \texttt{v}*,
\texttt{ω}* \\
N04 N06 N07 N08 & \texttt{eqsELei00} & \texttt{almightyWs} \\
N04 N06 N07 N14 & \texttt{My4ArrayComponents}, \texttt{TableDelta} &
\texttt{TableMetric} \\
N04 N10 N11 N12 & \texttt{eextract}, \texttt{makeAntiSym},
\texttt{makeSym}, \texttt{skelx} & \\
N06 N07 N11 N12 & \texttt{tau16}* & \\
N01 N02 N10 & \texttt{EA}, \texttt{hUSE} & \texttt{B}*, \texttt{J}*,
\texttt{g\_AB}*, \texttt{h}*, \texttt{unit} \\
N01 N05 N10 & \texttt{Over{[}σ22,\_{]}}, \texttt{anti4by4},
\texttt{symm4by4}, \texttt{txA}, \texttt{txA0}, \texttt{txAB},
\texttt{txABC}, \texttt{η16}, \texttt{η2244}, \texttt{σ22} &
\texttt{txABCD} \\
N01 N10 N12 & \texttt{coupled}, \texttt{coupled1}, \texttt{coupled2},
\texttt{linTrans}, \texttt{linTransAntiSymmetric},
\texttt{linTransSymmetric}, \texttt{negativeTrMM}, \texttt{positiveTrMM}
& \\
N02 N08 N09 & & \texttt{g}* \\
N04 N05 N10 & \texttt{MergeSetsStep}, \texttt{dualRawSets},
\texttt{showCoupledEquations} & \texttt{rawSets} \\
N04 N06 N07 & \texttt{ein} & \\
N04 N06 N14 & \texttt{g44} & \\
N04 N08 N14 & & \texttt{Lovelock1a}, \texttt{Lovelock2a} \\
N04 N10 N12 & & \texttt{blocks} \\
N04 N13 N14 & \texttt{test0} & \\
N10 N11 N12 & \texttt{Au}*, \texttt{Ax}*, \texttt{E\_((A))\^{}α}*,
\texttt{EinsteinG}, \texttt{LagrangianΨ16}*, \texttt{Lj},
\texttt{Lsquared}*, \texttt{RS}, \texttt{RicciΓ}, \texttt{RieΓ},
\texttt{SAB}, \texttt{SAB1}, \texttt{SAB2}, \texttt{T16\^{}α},
\texttt{T\^{}α}*, \texttt{Universe}*, \texttt{ass}, \texttt{con},
\texttt{constraintTrig}, \texttt{constraintVars}, \texttt{constraintX},
\texttt{covariantDiffMatrix}, \texttt{detgg}, \texttt{eAa},
\texttt{eAa04}, \texttt{eL}, \texttt{eLa}, \texttt{e\_((A))\^{}α},
\texttt{e\_α\^{}((A))}, \texttt{eηe}, \texttt{g4488},
\texttt{g\_(μν∣α)}*, \texttt{ginv}, \texttt{gtry},
\texttt{gtrye\_((A))\^{}α}, \texttt{gtrye\_α\^{}((A))}, \texttt{iid},
\texttt{makeSubstitution}, \texttt{qid}, \texttt{result}, \texttt{rt},
\texttt{sfψ16Aa}, \texttt{sg}, \texttt{sge\_((A))\^{}α},
\texttt{sge\_α\^{}((A))}, \texttt{sgtrye\_((A))\^{}α},
\texttt{sgtrye\_α\^{}((A))}, \texttt{spinCoefficients},
\texttt{spinCoeffs}, \texttt{sreplaceZ}, \texttt{ssgm4488},
\texttt{subsDefects}, \texttt{sx0x4}, \texttt{szt}, \texttt{sω},
\texttt{sωf}, \texttt{sωfgtry}, \texttt{sωgtry}, \texttt{useDSQRT},
\texttt{useT16}, \texttt{usegT16}, \texttt{usingLagrangianF16massive}*,
\texttt{Γ}, \texttt{Γ\_j}*, \texttt{Ψ16}, \texttt{Ψ16lower},
\texttt{Ψ16upper}, \texttt{Ω16}, \texttt{η99}, \texttt{ωmat},
\texttt{ωs}, \texttt{ωμIJ}, \texttt{ωω} & \texttt{La}, \texttt{Lg},
\texttt{TU\^{}μν}*, \texttt{eLazt}, \texttt{imes} \\
N01 N05 & \texttt{σ88} & \\
N01 N10 & \texttt{B1}*, \texttt{E\^{}A\_a}, \texttt{E\^{}a\_A},
\texttt{F\_A\^{}a}, \texttt{F\_a\^{}A}, \texttt{JT4},
\texttt{Jcomplement}, \texttt{JcomplementProjection8},
\texttt{Jcomplementcc}, \texttt{MS1S2}, \texttt{Myid}, \texttt{T4cc},
\texttt{T7cc}, \texttt{\textbackslash{}{[}DoubleStruckU{]}},
\texttt{antisymm16}, \texttt{base4by4}, \texttt{basisProjections16},
\texttt{basisProjections8}, \texttt{caΨccΨ}, \texttt{caΨccΨ2},
\texttt{eqsRQQ7Q8}, \texttt{evaluesAB}, \texttt{evecsAB},
\texttt{findComplexStucture}*, \texttt{gridx}, \texttt{jjcc},
\texttt{myid}, \texttt{parallel04},
\texttt{possibleComplexStructureIndex},
\texttt{possibleComplexStructureIndex2},
\texttt{possibleComplexStructureMatrix},
\texttt{possibleComplexT4StructureIndex}, \texttt{projection04},
\texttt{realProjection16}, \texttt{realProjection8}, \texttt{sJu},
\texttt{sQ7}, \texttt{sQ7Q8}, \texttt{sQ8}, \texttt{sRQQ7Q8},
\texttt{seqsRQQ7Q8}, \texttt{sixers}, \texttt{sixersIndices},
\texttt{sixersJ}, \texttt{type1SpinorBasisMatrix},
\texttt{type2SpinorBasisMatrix}, \texttt{unit3}, \texttt{xxx},
\texttt{yyy}, \texttt{Ψ16cc}, \texttt{Ψ16ccA}, \texttt{Ψ16cccc},
\texttt{τABC} & \\
N02 N07 & \texttt{g71} & \\
N02 N09 & & \texttt{α}*, \texttt{β}*, \texttt{δ}*, \texttt{ε}*,
\texttt{η}*, \texttt{κ}*, \texttt{λ}*, \texttt{μ}*, \texttt{ν}*,
\texttt{ξ}*, \texttt{ρ}*, \texttt{χ}* \\
N04 N10 & \texttt{blockPartition} & \\
N04 N14 & \texttt{gg44} & \texttt{makeSym00} \\
N05 N10 & \texttt{S44αβ}, \texttt{g3}, \texttt{g\_αβ}, \texttt{sε},
\texttt{γABs} & \texttt{η44} \\
N06 N07 & \texttt{MatrixMetric71}, \texttt{a4} & \\
N08 N14 & & \texttt{Lovelock3a} \\
N10 N12 & \texttt{ConvertMapleToMathematica}, \texttt{GetTokens},
\texttt{IsAlpha}, \texttt{IsDigit}, \texttt{IsSpace},
\texttt{MapleToMathematicaFunction}, \texttt{ParseTokens},
\texttt{ToMathematicaString}, \texttt{helpf},
\texttt{mapleztStringToMathematica}, \texttt{snewfψ16Aa},
\texttt{sreplacenZ}, \texttt{sta}, \texttt{sti}, \texttt{times},
\texttt{ω\_((B)(A)\ )}* & \texttt{ProcessASTNode} \\
N11 N12 & \texttt{LagrangianF16massive}* & \texttt{matTest},
\texttt{resolveWithRemainder} \\
\end{longtable}

\end{document}
